% =====================================================================
%  Manuscript for JGR: Planets.
%
%  Title: Difference of Lunar Regolith Thermal Conductivity K_d at the
%         Apollo 15 and 17 Heat-Flow Boreholes.
%
%  Build:  pdflatex letter ; bibtex letter ; pdflatex letter ; pdflatex letter
%
%  CHANGE TRACKING
%  ---------------
%  ...  inserted / revised text  (green)
%    removed text, struck through for the record  (red)
%   author query / to-do  (orange)
%  DISPLAY MODE: set \def\displaymode{...} below to one of
%    review  | newonly | clean   (see the block before the title).
%    review  = full markup;  newonly = only the new wording, highlighted;
%    clean   = submission PDF (no markup at all).
% =====================================================================
\documentclass[11pt]{article}

\usepackage[a4paper,margin=2.2cm]{geometry}
\usepackage[T1]{fontenc}
\usepackage{lmodern}
\usepackage{microtype}
\usepackage{amsmath,amssymb}
\usepackage{siunitx}
\sisetup{detect-all,per-mode=symbol}
% `year' clashes with the TeX \year primitive, which makes
% \si{...\per\year} throw "Missing number". Declare an explicit unit.
\DeclareSIUnit{\year}{yr}
\usepackage{graphicx}
\usepackage{booktabs}
\usepackage{caption}
\captionsetup{font=small,labelfont=bf,labelsep=period}
\usepackage{xcolor}
\usepackage[hidelinks,colorlinks=true,linkcolor=blue!50!black,
            citecolor=blue!50!black,urlcolor=blue!60!black]{hyperref}
\usepackage[round,authoryear]{natbib}
\bibliographystyle{plainnat}
\setlength{\bibsep}{2pt}
\renewcommand{\arraystretch}{1.15}

% --- JGR/AGU review: continuous line numbers in the left margin --------------
\usepackage[mathlines]{lineno}

% --- Float placement: keep figures near their references --------------------
\usepackage{placeins}             % \FloatBarrier
\renewcommand{\topfraction}{0.95}
\renewcommand{\bottomfraction}{0.85}
\renewcommand{\textfraction}{0.05}
\renewcommand{\floatpagefraction}{0.85}
\setcounter{topnumber}{2}
\setcounter{bottomnumber}{2}
\setcounter{totalnumber}{4}

% --- Tighten whitespace around floats ---------------------------------------
\setlength{\floatsep}{6pt plus 2pt minus 2pt}
\setlength{\textfloatsep}{8pt plus 2pt minus 4pt}
\setlength{\intextsep}{6pt plus 2pt minus 2pt}
\setlength{\abovecaptionskip}{4pt}
\setlength{\belowcaptionskip}{2pt}

% --- Tighten equation vertical spacing --------------------------------------
\setlength{\abovedisplayskip}{6pt plus 2pt minus 2pt}
\setlength{\belowdisplayskip}{6pt plus 2pt minus 2pt}
\setlength{\abovedisplayshortskip}{3pt plus 1pt}
\setlength{\belowdisplayshortskip}{3pt plus 1pt}

% --- Tighten section/paragraph heading spacing ------------------------------
\usepackage{titlesec}
\titlespacing*{\section}{0pt}{10pt plus 2pt minus 2pt}{4pt plus 1pt}
\titlespacing*{\subsection}{0pt}{8pt plus 2pt minus 2pt}{3pt plus 1pt}
\titlespacing*{\paragraph}{0pt}{5pt plus 1pt minus 1pt}{1em}

% --- Plain-language summary box (drop when moving to the AGU class) ----------
\newenvironment{plainlanguagesummary}{%
  \par\medskip\noindent\textbf{Plain Language Summary.}\itshape\small}{\par\medskip}

% --- Convenience macros -----------------------------------------------------
\newcommand{\Hayne}{Hayne et al.\ (\citeyear{hayne2017})}
\newcommand{\Vas}{Vasavada \emph{et~al.} (\citeyear{vasavada2012})}
\newcommand{\dT}{\Delta T}
\newcommand{\zborestem}{80\,\si{\centi\meter}}

% =====================================================================
%  Change-tracking macros (replace the latexdiff \DIFadd/\DIFdel machinery)
% =====================================================================
\usepackage[normalem]{ulem}
% let \sout survive inline math inside 
\makeatletter
\let\sout@orig\sout
\renewcommand{\sout}[1]{\ifmmode\text{\sout@orig{\ensuremath{#1}}}\else\sout@orig{#1}\fi}
\makeatother

% =====================================================================
%  DISPLAY MODE -- change the ONE word below to switch how the document
%  renders. No other edits needed.
%
%    review   : full markup -- new text green, old text struck through in
%               red, author notes shown in orange. (Use while editing.)
%    newonly  : show ONLY the new wording -- new text highlighted green,
%               the struck-out old text and the notes HIDDEN. (Read the
%               revised paper with the additions still visible.)
%    clean    : final submission look -- new text plain black, old text
%               and notes removed entirely, nothing highlighted.
%
%  ONE-CLICK PRINTABLE VERSION: compile letter_clean.tex instead of this
%  file (or run ./build_clean.sh) to force `clean` mode and get a
%  flattened, markup-free PDF ready to print/submit.
% =====================================================================
% Final draft: all tracked changes accepted and flattened into the text.
% These no-op definitions remain only so the file compiles if a stray
% \new/\old/\note ever reappears from a copy-paste; they render plainly.
\newcommand{\new}[1]{#1}
\newcommand{\old}[1]{}
\newcommand{\note}[1]{}

% --- Title ------------------------------------------------------------------
% Alt title kept for reference:
% Estimates of Lunar Regolith Thermal Conductivity K_d at the Apollo 15 and 17 Heat-Flow Boreholes
\title{\bfseries  Difference of Lunar Regolith Thermal Conductivity $K_d$ at the Apollo 15 and 17 Heat-Flow Boreholes}


% --- Author block -----------------------------------------------------------
% TODO before submission: confirm co-author first initials and Larsson's
% affiliation; on switching to the AGU class, replace this whole block with
% \authors{} + \affiliation{} + \correspondingauthor{}.
\author{%
  Ramon~III~P.~Gregorio$^{1,2,3}$,\quad
  Richard~Larsson$^{1}$,\quad
  Takayoshi~Yamada$^{2}$,\quad
  Yasuko~Kasai$^{1,2}$%
}

% --- Manual affiliation footer (the AGU class will absorb this) -------------
\makeatletter
\renewcommand\@maketitle{%
  \newpage \null \vskip 2em%
  \begin{center}%
    {\LARGE \@title \par}%
    \vskip 1em%
    {\large \lineskip .5em \@author \par}%
    \vskip .5em%
    {\small
      $^{1}$Institute of Science Tokyo, Tokyo, Japan.\\
      $^{2}$National Institute of Information and Communications
      Technology (NICT), Tokyo, Japan.%
      \\
     $^{3}$Ateneo de Davao University, Davao, Philippines.\\
      Correspondence to: \texttt{ Yasuko Kasai, kasai.y.aa@m.titech.ac.jp}.%
    }%
    \vskip 1em {\@date}%
  \end{center}\par \vskip 1.5em%
}
\makeatother

\date{\today}

\begin{document}
\linenumbers
\maketitle


% NOTE 2026-07-04: AGU accepts FREE-FORMAT initial submissions; the AGUTeX
% port below is required only at acceptance. The template zip is no longer
% at the old agu.org URL -- fetch it from the journal's Author Resources
% page (or Overleaf's official AGU template) when porting.
% TODO before submission: (i) port to the AGUTeX (agujournal2019)
% class; (ii) add the AGU-required CRediT author-contribution
% statement (needs each co-author's roles); (iii) add the funding
% statement (grant numbers from Kasai Lab / NICT, or "no external
% funding").

% =====================================================================
\begin{abstract}
\noindent
Apollo 15 and 17 are the only sites on the Moon where subsurface
regolith temperatures have been measured in situ. The restored
1971--1977 Heat-Flow Experiment (HFE) record is the sole direct
constraint on the meter-scale regolith thermal conductivity $K_d$, yet
current thermal models apply a single Moon-wide value
\citep{hayne2017,martinez2021}. We retrieve $K_d$ separately at each
borehole from the restored record \citep{nagihara2018} by sweeping it
against the meter-scale-sensor RMSE (80--234 cm) under the
\citet{hayne2017} form, excluding borestem-contaminated shallow
sensors. Conditional on the published basal heat fluxes ($21$ and $16$
mW m$^{-2}$; \citealp{langseth1976}), the retrieval yields
$K_{d,\mathrm{A15}}^{*}=4.60^{+0.86}_{-0.21}$ and
$K_{d,\mathrm{A17}}^{*}=7.08^{+0.53}_{-0.52}$ mW m$^{-1}$ K$^{-1}$
($1\sigma$ bootstrap), reducing the sensor RMSE by 8\% at Apollo 15 and
55\% at Apollo 17 relative to the global $K_d=3.4$. The inter-site
contrast, $\Delta K_d^{*}=2.3^{+0.7}_{-0.9}$ mW m$^{-1}$ K$^{-1}$ (95\%
CI $[-0.1,\,3.6]$), marginally includes zero
($P_{\rm boot}\approx 0.03$); the ordering
$K_{d,\mathrm{A17}}^{*}\!>\!K_{d,\mathrm{A15}}^{*}$ nonetheless survives
the basal-flux reanalysis ($\ge$99\% across the published $Q_b$
envelope), though the absolute values do not. Both retrievals lie
within porosity, composition, and orbital benchmarks and supply the
meter-scale boundary condition needed by subsurface radiative-transfer
retrievals such as the upcoming TSUKIMI terahertz mission.
\end{abstract}

\begin{plainlanguagesummary}
The Moon's interior is insulated from the daily surface heating by the
compacted regolith (broken rock and dust) in roughly the top meter. The
numbers describing that insulation come from formulas fit to satellite
measurements of the surface temperature, not to any direct sub-surface
data. Sub-surface temperatures were measured at only two locations --- by
thermometers the Apollo 15 and 17 astronauts emplaced up to 1.4 and
2.3 m deep in 1971--1972 --- and the complete records were recently
recovered from the original mission tapes. We use those records to
measure, at each borehole, the regolith conductivity about a meter down,
where the daily temperature swing has died away and the grain-to-grain
part of the conductivity has reached its deep, constant value. The
Hayne (2017) global value misses these temperatures by about 1.1 K at
Apollo 15 and 0.9 K at Apollo 17; refitting each site separately halves
the Apollo 17 mismatch, and the two sites disagree on the conductivity
by roughly a factor of 1.5. These site-specific values provide the
sub-surface temperature input that upcoming lunar mapping missions need
to interpret what their instruments see through the regolith.
\end{plainlanguagesummary}

\medskip
\noindent\textbf{Key Points}
\begin{itemize}\itemsep2pt
  \item Meter-scale regolith conductivity
        $K_d$ is retrieved per site from the restored Apollo 15 and 17
        heat-flow records (80--234 cm).
  \item Apollo 15 gives
        $K_d=4.60^{+0.86}_{-0.21}$ and Apollo 17 gives
        $7.08^{+0.53}_{-0.52}$ mW m$^{-1}$ K$^{-1}$ ($1\sigma$
        bootstrap).
  \item Both sites reject the Hayne et al.\ (2017) global $K_d{=}3.4$
        at 95\%; the contrast is marginal ($p{\approx}0.03$) but its
        ordering is robust ($\ge$99\%).
\end{itemize}


% =====================================================================
\section{Introduction}\label{sec:intro}

Lunar sub-surface thermal modeling presently relies on globally averaged
conductivity parameterizations. We use the term
\emph{meter-scale regolith} throughout to mean
the subsurface layer constrained by the
Apollo HFE temperature sensors used in this retrieval---specifically,
80--139 cm for Apollo 15 and 130--234 cm for Apollo 17. At these depths the
daily and annual thermal skins ($\lesssim 10$ cm and $\sim 50$ cm
respectively) have damped to negligible amplitude, and the conductivity in
both standard parameterizations has
converged on its depth-independent asymptote
$K_d$; under the Hayne (2017) form, the same $K_d$ value
extrapolates to all greater depths within the regolith column. The meter-scale regolith conductivity is supplied by formulations
calibrated against orbital Diviner brightness temperatures: the
smooth-exponential $K(T,z)$ of \citet{hayne2017}, in which a single
meter-scale asymptote $K_d$ is fit to a global
average and applied moon-wide, and the temperature- and density-dependent
$K(T,\rho)$ of \citet{martinez2021}, which extends the same functional
family with additional physical content but is itself calibrated globally
and applied uniformly. Both parameterizations probe only the diurnal skin.
The extent to which the resulting global $K_d$ represents any individual
landing site, and the magnitude of any site-to-site departure, cannot be
determined from orbital data alone.

Direct in-situ data exist only from the Apollo Heat-Flow Experiment (HFE).
At Apollo 15 and 17, fiberglass borestems emplaced by the astronauts
carried thermometers to depths of up to 1.4 m (Apollo 15) and 2.3 m
(Apollo 17) \citep{langseth1972,langseth1976}; these two boreholes remain the only
lunar locations at which a sub-surface temperature profile has been
measured (Fig.~\ref{fig:intro-probe}). The restored 1971--1977 HFE record
\citep{nagihara2018} provides the data used here. Each borehole also
carries a published basal heat flux \citep{langseth1976,nagihara2018}
which, together with the local $K_d$, sets the
meter-scale temperature gradient. The site-specific inputs
(latitudes, albedos, emissivities, basal fluxes) are collected with their
original sources in Table~\ref{tab:params}.

We retrieve $K_d$ separately at each site by holding the \citet{hayne2017}
functional form fixed and sweeping $K_d$ alone against the
meter-scale-sensor RMSE. The retrieval is placed in
context by two additional forward evaluations against the same
meter-scale-sensor set: the \citet{hayne2017} form
at its published global $K_d$, and the \citet{martinez2021} $K(T,\rho)$
model at its published coefficients (Table~\ref{tab:headline}). Neither
published model is refit here, so the \citet{martinez2021} comparison
quantifies the performance of the published global $K(T,\rho)$ at each
Apollo site without per-site freedom. A full per-site retrieval under the
\citet{martinez2021} form requires a different choice of free parameter from
$K_d$ and is deferred to a follow-on study (\S\ref{sec:disc-prior}).

The per-site $K_d^{*}$ values matter for any sub-surface radiative-transfer
(RTM) retrieval that observes the Moon at wavelengths penetrating below the
diurnal skin: future microwave
and terahertz sub-surface
sounders all need a meter-scale temperature column $T(z)$
as a forward-model boundary condition that orbital instruments cannot
measure directly. One concrete example is the upcoming TSUKIMI mission
(Lunar Terahertz Surveyor for Kilometer-scale Mapping,
\citealp{tsukimi2022}), whose forward model requires the sub-diurnal-skin temperature column
that only the in-situ $K_d$ can anchor.

Below the diurnal skin the modeled column is controlled by
$K_d$ through both the rectified lunation mean and the steady gradient
$Q_b/K_d$, so the $T(z)$ ingested by any such retrieval inherits its
$K_d$ dependence; a
factor-of-two error in $K_d$ propagates through the
meter-scale temperature profile and into any retrieval
that uses it. The values reported here, evaluated in the \citet{hayne2017}
thermal column, supply that boundary condition from the only in-situ data
available.\label{sec:intro-tsukimi}

The retrieval is driven by an
idealized synodic insolation cycle evaluated at each site's latitude
(\S\ref{sec:methods}); because only stability-window mean temperatures
enter the fit, sub-cycle timing plays no role, and the neglected
ephemeris terms (solar declination, orbital eccentricity) perturb the
cycle-mean insolation at or below the percent level
(\S\ref{sec:methods}). The remaining sections present the
data and methods (\S\ref{sec:methods}), the per-site retrieval with the
headline RMSE table (\S\ref{sec:results}), the interpretation with its
uncertainty limits (\S\ref{sec:discussion}), and the conclusions.

\begin{figure}[htbp]
\centering
\includegraphics[width=\linewidth]{figures/fig_intro_probe.pdf}
\caption{Apollo HFE measurement geometry and the two published $K(z)$ forms
tested. \textbf{(a)} Borestem emplacements at both sites, to a shared depth
scale: markers are the real sensor depths, filled where a sensor enters the
retrieval ($z\ge 80$ cm) and open in the borestem-disturbed zone
($z<80$ cm, coral) that is excluded; the two probe strings per site are
offset left/right. Apollo~17 reaches 2.3 m (16 retained sensors), Apollo~15
1.4 m (7). \textbf{(b)} Conductivity profiles $K(z)$ at $T=250$ K:
\citet{hayne2017} at $K_d^{\mathrm{publ.}}=3.4$ mW m$^{-1}$ K$^{-1}$ (solid
teal) and \citet{martinez2021} $K(T,\rho)$ (dashed violet), sharing the same
Hayne-form $\rho(z)$.}
\label{fig:intro-probe}
\end{figure}

\begin{figure}[htbp]
\centering
\includegraphics[width=\linewidth]{figures/fig_context_map.pdf}
\caption{Apollo Heat-Flow Experiment landing-site context. \textbf{(a)}
Global LOLA topography and \textbf{(b)} modeled mean surface temperature,
each with the two boreholes marked; dashed rectangles on (a) show the zoom
regions. \textbf{(c,d)} LOLA shaded-relief zooms ($\pm 15^\circ$ longitude,
$\pm 10^\circ$ latitude) resolving the Imbrium--Apennine mare/highland margin
at Apollo~15 and the Taurus--Littrow highlands at Apollo~17
\citep{korotev2005,wieczorek2006,crawford2015}. Borehole coordinates from
the ALSEP gazetteer.}
\label{fig:context-map}
\end{figure}

% =====================================================================
\section{Data and Methods}\label{sec:methods}

\subsection{Apollo HFE meter-scale-sensor temperatures}
The retrieval requires one equilibrium temperature $T_\mathrm{eq}$ per
sensor. The restored 1971--1977 record \citep{nagihara2018} provides this
quantity once the early portion of each record, which contains residual
transients from drilling and emplacement, documented mission-operations
disturbances \citep{nagihara2018}, and slow instrumental drift, is
excluded. For each sensor we extract automatically the longest trailing
segment that is statistically indistinguishable from a flat steady state
and average within it.

For each sensor, candidate start points covering 55--80\%
of the record
(13 equally-spaced candidates; those retaining less than 20\% of the
record are excluded by a minimum-tail guard, capping the effective range
at 80\%) are evaluated by ordinary-least-squares linear
fit to all temperatures between the candidate point and the end of the
record. The accepted \emph{stability window} is the earliest candidate for
which the linear-fit slope satisfies
$|d\langle T\rangle/dt|<0.08\,\si{\kelvin\per\year}$
($\approx 2.2\times 10^{-4}\,\si{\kelvin\per\day}$): the longest tail
consistent with thermal equilibrium at this stringency. The adopted
threshold is approximately $4.6\times$ stricter than the
$10^{-3}\,\si{\kelvin\per\day}$ value sometimes used in HFE analyses, and
was chosen to suppress bias from residual post-disturbance drift. If no
candidate satisfies the criterion, the last 25\% of the record is used as a
fallback. The strict criterion is met by 3 of 7 retained sensors at
A15 and 6 of 16 at A17; the remainder take the fallback window, with
residual trailing slopes of 0.11--0.47 K yr$^{-1}$. The threshold sweep
below bounds the sensitivity to the choice of criterion, but
residual common-epoch drift itself is not propagated as a systematic, and
the within-window scatter reported as the per-sensor uncertainty
characterizes short-term stability only --- it does not weight the
headline RMSE or the bootstrap (\S\ref{sec:limitations}). We report the window-mean equilibrium temperature
$T_\mathrm{eq}=\langle T(z)\rangle_{\Delta t}$ together with its
within-window standard deviation, which serves as the per-sensor
measurement uncertainty in downstream statistics. Sensors with central
depth $z<\zborestem$ (the borestem-contaminated zone;
\S\ref{sec:borestem}) are plotted but excluded from the RMSE.

This construction should not be read as a claim that every archived
Apollo trace had reached exact zero-drift steady state. It is an estimator
for the undisturbed periodic-equilibrium temperature of the regolith after
the borehole disturbance has vanished. In the nominal reduction the strict
flat-trend criterion is satisfied by 3/7 Apollo~15 and 6/16 Apollo~17
meter-scale sensors; the remaining retained sensors use the late-record
fallback tail and therefore carry residual post-emplacement drift as a
systematic component of $T_\mathrm{eq}$. The forward model below is thus
compared to equilibrium-temperature estimates, not to raw
instantaneous measurements.

Figure~\ref{fig:timeline} shows the full per-probe temperature record at
each site, with the documented disturbance intervals shaded
\citep{nagihara2018}, the selected stability windows marked, and the fitted
long-term drift $d\langle T\rangle/dt$ reported per sensor.

The retrieval is insensitive to the precise threshold. Sweeping the
threshold over $0.04$--$\SI{0.16}{\kelvin\per\year}$ (a factor of four
around the adopted value) shifts $K_d^{*}$ by at most
$0.02\,\si{\milli\watt\per\meter\per\kelvin}$ at A15
(0.4\% of nominal) and
$0.05\,\si{\milli\watt\per\meter\per\kelvin}$ at A17
(0.6\%), in both cases
an order of magnitude below the bootstrap $1\sigma$ uncertainty
(\S\ref{sec:results-kd}). The adopted threshold therefore lies within the
flat band of the threshold response.


\begin{figure}[!htbp]
\centering
\includegraphics[width=\linewidth]{figures/fig_apollo_timeline_probes.pdf}
\caption{Per-probe stability-window timeline at Apollo~15 and~17 (4
probes). Each probe panel shows, above, the raw HFE traces with coral bands
marking documented data-quality events and the selected stability windows
outlined; and, below, a depth-coded per-sensor Gantt strip broken at
outages. Selection criteria and the role of the window-mean
$T_\mathrm{eq}$ are described in \S\ref{sec:methods}.}
\label{fig:timeline}
\end{figure}

\subsection{Surface solar geometry and incident insolation}
The solver is forced
with an idealized synodic insolation cycle: the daytime solar zenith
angle at each site is approximated by
$\cos\theta_\odot(t)=\cos\phi\,\cos(2\pi t/P)$, with $\phi$ the site
latitude and $P=29.531$ d the mean synodic period. Because the retrieval
compares stability-window mean temperatures (\S\ref{sec:methods})
rather than time-resolved phases, ephemeris-level timing corrections do
not enter the fit. The neglected solar declination ($\pm1.54^\circ$) and
orbital eccentricity ($\pm3.3\%$ peak-to-peak in instantaneous flux,
$<0.1\%$ in the cycle mean) perturb the mean insolation at or below the
percent level, comparable to and subsumed by the Bond-albedo sensitivity
of Table~\ref{tab:errorbudget}. The Moon has no atmosphere, so
the solar irradiance incident on the local horizontal surface is
\begin{equation}
F_\odot(t) = \max\bigl(0,\ S_0\cos\theta_\odot(t)\bigr),
\label{eq:insolation}
\end{equation}
with $S_0 = 1361\,\si{\watt\per\meter\squared}$ \citep{kopp2011}.


\subsection{1-D solver and \Hayne\ thermophysics}\label{sec:solver}

The depth coordinate $z$ is defined positive downward, with the regolith
surface at $z=0$ and the lower boundary of the modeled column at
$z=z_\mathrm{max}=5$ m. Conservation of energy in one dimension gives the
heat equation
\begin{equation}
\rho(z)\,c_p(T)\,\frac{\partial T}{\partial t}
   = \frac{\partial}{\partial z}\!\Bigl[\,K(T,z)\,\frac{\partial T}{\partial z}\Bigr].
\label{eq:heat}
\end{equation}
At the upper boundary the skin temperature $T_s$ satisfies radiative
balance between absorbed insolation, emitted thermal radiation, and the
conductive flux into the column,
\begin{equation}
(1-A)\,F_\odot(t)\;-\;\varepsilon\sigma T_s^4
\;-\;\bigl(-K\,\partial_z T\bigr)_{\!z=0} \;=\; 0,
\label{eq:bc-top}
\end{equation}
with $\sigma=5.6704\times10^{-8}$ W m$^{-2}$ K$^{-4}$ the Stefan--Boltzmann
constant; equation~(\ref{eq:bc-top}) is solved iteratively
for $T_s$ at each time step. At the lower boundary a constant geothermal
flux is imposed,
\begin{equation}
K\,\partial_z T\big|_{z_\mathrm{max}} = Q_b,
\label{eq:bc-bot}
\end{equation}
with $z$ positive downward: the upward geothermal flux maintains a
positive (warming-with-depth) temperature gradient, matching the solver's
basal ghost condition. The PDE is discretized on a geometric grid ($\Delta z_0=2$ mm, growth ratio
1.08) and integrated with a Crank--Nicolson scheme at
$\Delta t = \SI{1800}{\second}$ on a lunation-commensurate time grid whose
final (cycle-closing) step is marched explicitly, with material properties
re-evaluated at the time-step midpoint by two corrector sweeps
(the implicit-midpoint prescription; both details are required for the
periodic state to converge at second order in $\Delta t$).
The periodic steady state is computed by a flux-anchored outer
iteration rather than a fixed-length spin-up: short multi-lunation
integrations converge the diurnal skin, after which the sub-skin
mean profile is reconstructed from the steady-state mean-flux closure
$d\langle T\rangle/dz=\bigl[Q_b-u_{\rm rect}(z)\bigr]/K(\langle
T\rangle,z)$, where the rectified (eddy-correlation) flux
$u_{\rm rect}=\langle K\,\partial_z T\rangle-K(\langle
T\rangle)\,\partial_z\langle T\rangle$ is measured from the resolved
cycle; the two steps are alternated until the anchor temperature at
$z=0.55$ m is stable to $5\times10^{-3}$ K. The skin integration is
confined to a shallow sub-column ($z\le 0.70$ m) with $Q_b$ imposed at
its base, while the deep column is supplied entirely by the closure
reconstruction and never time-stepped; this decouples the spin-up from
the $\sim$$10^3$-lunation relaxation of the full 5-m domain, so a modest
inner length is free of the initial condition. Convergence is judged on
the retrieved $K_d^{*}$ itself: it is stable to
$<\!5\times10^{-3}$ mW m$^{-1}$ K$^{-1}$ once $n_\mathrm{inner}\gtrsim 64$
at Apollo 17 and $\gtrsim 96$ at the lower-conductivity, slower-relaxing
Apollo 15 site, which sets the production $n_\mathrm{inner}=96$; the
truncated and full-column constructions agree on $K_d^{*}$ to
$2\times10^{-3}$ mW m$^{-1}$ K$^{-1}$, confirming the truncation is
unbiased. Initial temperature guesses of 240 and 260 K yield mean
profiles agreeing to $\le 0.03$ K everywhere in the column, and
grid-and-timestep refinement, anchor-depth variation over
$0.45$--$0.70$ m, and the $K_d$-grid spacing bound the residual solver
systematic on $K_d^{*}$ at $\lesssim 0.04$ mW m$^{-1}$ K$^{-1}$ --- well
below the bootstrap uncertainty (\S\ref{sec:results-kd}).
Temporal convergence is certified in the reported quantity itself:
halving $\Delta t$ from 1800 to 900 to \SI{450}{\second} moves
$K_d^{*}$ by $\le 0.018$ (A17) and $\le 0.001$ (A15)
mW m$^{-1}$ K$^{-1}$, and the Richardson extrapolation of an
intentionally first-order variant of the scheme converges to the same
$K_d^{*}$(A17) to four significant digits --- two independent
discretizations agreeing on the continuum limit. The solver is further
verified against an independent line-for-line C++ implementation, with
which it agrees to $10^{-12}$ K over multi-lunation integrations; all
certification artifacts are archived with the pipeline
(\S\ref{sec:openresearch}). The
\citet{hayne2017} thermophysical inputs are adopted as given in Appendix A
of the original paper:
\begin{align}
K(T,z) &= \Bigl\{K_s + (K_d-K_s)\bigl[1-e^{-z/H}\bigr]\Bigr\}
         \cdot\Bigl[1 + \chi\bigl(T/T_\mathrm{ref}\bigr)^3\Bigr], \label{eq:K} \\
\rho(z) &= \rho_d - (\rho_d-\rho_s)\,e^{-z/H}, \label{eq:rho}\\
c_p(T)  &= c_0 + c_1 T + c_2 T^2 + c_3 T^3 + c_4 T^4. \label{eq:cp}
\end{align}
Throughout, $K_d$ denotes the deep asymptote of the \emph{contact}
component $K_c(z)$ of equation~(\ref{eq:K}); the total conductivity
$K(T,z)$ carries in addition the radiative factor
$1+\chi(T/T_\mathrm{ref})^3$, which at the retained-sensor temperatures
($\approx$252--257 K) multiplies the contact term by $\approx$2.0. At the
retrieved $K_d^{*}$ the mean \emph{effective} conductivity at the sensors
is 9.3 (A15) and 14.6 (A17) mW m$^{-1}$ K$^{-1}$; literature comparisons
must respect this distinction (\S\ref{sec:disc-prior}).
Numerical values and sources for every coefficient, together with the
site-specific measured quantities (latitude, albedo, emissivity, basal heat
flux), are listed in Table~\ref{tab:params}. The basal heat flux
reanalyses of \citet{saito2007,nagihara2018} are folded into the
sensitivity sweep of \S\ref{sec:results-kd}.

\begin{table}[!htbp]
\centering
\footnotesize
\setlength{\tabcolsep}{4pt}
\renewcommand{\arraystretch}{1.10}
\caption{\textbf{Fixed inputs to the 1-D heat solver.} Three blocks:
(A)~the \citet{hayne2017} conductivity-form coefficients (eq.~\ref{eq:K}),
with $K_d$ the single parameter retrieved in \S\ref{sec:methods-kd};
(B)~the four published coefficients of the \citet{martinez2021} $K(T,\rho)$
comparison model (eq.~\ref{eq:K-ms}, forward only in the headline table;
see also \S\ref{sec:disc-prior}); (C)~site-specific measured quantities
shared by both conductivity models. The shared $\rho(z)$ uses the
Hayne-form profile (eq.~\ref{eq:rho}); the specific heat is the
\citet{hemingway1973} polynomial
$c_p(T)=c_0+c_1 T+c_2 T^2+c_3 T^3+c_4 T^4$ with
$(c_0,\ldots,c_4)=(-3.61,\,2.74,\,2.36\!\times\!10^{-3},\,
{-1.23}\!\times\!10^{-5},\,8.91\!\times\!10^{-9})$.}
\label{tab:params}
\begin{tabular}{@{} l l l p{0.20\linewidth} p{0.22\linewidth} @{}}
\toprule
Parameter & Symbol & Value & Unit & Source \\
\midrule
\multicolumn{5}{@{}l}{\textit{(A) Hayne et al.\ (2017) smooth-exponential conductivity, all values from \citet{hayne2017} App.~A unless noted}} \\[2pt]
Surface conductivity        & $K_s$              & $7.4\times10^{-4}$  & W\,m$^{-1}$\,K$^{-1}$ &  \\
Deep conductivity (global)  & $K_d$ (publ.)      & $3.4\times10^{-3}$  & W\,m$^{-1}$\,K$^{-1}$ &  \\
Deep conductivity (per-site)& $K_d^{*}$          & \S\ref{sec:results-kd} & W\,m$^{-1}$\,K$^{-1}$ & \textbf{this work} \\
$K(z)$ transition scale     & $H$                & 0.06                & m                     &  \\
Radiative coefficient       & $\chi$             & 2.7                 & ---                   &  \\
Radiative normalization     & $T_\mathrm{ref}$   & 350                 & K                     &  \\
\midrule
\multicolumn{5}{@{}l}{\textit{(B) Mart\'inez \& Siegler (2021) $K(T,\rho)$ \emph{(forward only)}, coefficients from \citet{martinez2021}}} \\[2pt]
Contact, $\rho$-slope       & $A_1$              & $5.082\times10^{-6}$   & W\,m$^{-1}$\,K$^{-1}$ (kg\,m$^{-3})^{-1}$ &  \\
Contact, intercept          & $A_2$              & $-5.10\times10^{-3}$   & W\,m$^{-1}$\,K$^{-1}$  &  \\
Radiative, $\rho$-slope     & $B_1$              & $2.002\times10^{-13}$  & W\,m$^{-1}$\,K$^{-4}$ (kg\,m$^{-3})^{-1}$ &  \\
Radiative, intercept        & $B_2$              & $-1.953\times10^{-10}$ & W\,m$^{-1}$\,K$^{-4}$  &  \\
\midrule
\multicolumn{5}{@{}l}{\textit{Shared by both conductivity models}} \\[2pt]
Surface bulk density        & $\rho_s$           & 1100                & kg\,m$^{-3}$           & \citet{hayne2017} \\
Deep bulk density           & $\rho_d$           & 1800                & kg\,m$^{-3}$           & \citet{hayne2017} \\
Specific heat               & $c_p(T)$           & polyn.\ (caption)   & J\,kg$^{-1}$\,K$^{-1}$ & \citet{hemingway1973} \\
\midrule
\multicolumn{5}{@{}l}{\textit{(C) Site-specific measured inputs}\quad\textit{(A15\,/\,A17)}} \\[2pt]
Latitude                    & $\phi$             & $26.1^\circ$\,/\,$20.2^\circ$\,N & --- & ALSEP gazetteer \\
Bond albedo                 & $A$                & $0.131$\,/\,$0.137$ & ---                   & \citet{vasavada2012} \\
Broadband emissivity        & $\varepsilon$      & $0.95$              & ---                   & \citet{bandfield2015} \\
Basal heat flux             & $Q_b$              & $21$\,/\,$16$       & mW\,m$^{-2}$          & \citet{langseth1976,nagihara2018} \\
Solar constant              & $S_0$              & $1361$              & W\,m$^{-2}$           & \citet{kopp2011} \\
\bottomrule
\end{tabular}
\end{table}

As an independent comparison we also evaluate the \citet{martinez2021}
regolith conductivity model. Where \citet{hayne2017} prescribes a
depth-dependent envelope through
$K(T,z)=[K_s+(K_d-K_s)(1-e^{-z/H})](1+\chi(T/T_\mathrm{ref})^3)$,
\citet{martinez2021} replaces that envelope with an explicitly
density-dependent two-mechanism decomposition,
\begin{equation}
K(T,\rho) = \underbrace{(A_1\rho + A_2)\,k_\mathrm{am}(T)}_{\text{contact (phonon)}}
          \;+\; \underbrace{(B_1\rho + B_2)\,T^3}_{\text{radiative}}.
\label{eq:K-ms}
\end{equation}
The first term is the contact (phonon) conductivity transmitted through the
network of grain-to-grain contacts; it carries the temperature dependence
of an amorphous solid through the low-temperature polynomial
$k_\mathrm{am}(T)$ of \citet{martinez2021} and scales linearly with the
local bulk density, reflecting the increased number of grain-to-grain
contacts per unit volume at higher packing. The second term is the
radiative contribution from infrared photon exchange across pore voids; it
scales as $T^3$ (the derivative of the Stefan--Boltzmann emission with
respect to temperature, characteristic of photon-mediated transport) and is
density-modulated with the opposite physical content, since denser packing
reduces the void fraction available for radiative exchange. The four
coefficients $(A_1,A_2,B_1,B_2)$ are calibrated globally against Diviner
brightness temperatures and are listed in Table~\ref{tab:params}. The
density profile $\rho(z)$ used in this comparison is the same Hayne-form
exponential used in the Hayne column, so the two models differ in their $K$
formulation alone.

The \citet{martinez2021} formulation has no single $K_d$ analogous to the
parameter retrieved under the Hayne form: both terms in
equation~(\ref{eq:K-ms}) vary continuously with depth through $\rho(z)$ and
with temperature through $T(z)$. The published model contains no free
parameter, so its evaluation at each site yields a single forward
prediction (Table~\ref{tab:headline}). To diagnose, in the discussion
(\S\ref{sec:disc-prior}), whether the published \citet{martinez2021} form
could fit both sites if it were granted one per-site lever analogous to the
Hayne $K_d$, we additionally perform an alternative retrieval in which the
asymptote of the Hayne-form density profile entering
equation~(\ref{eq:K-ms}) is rescaled by a single per-site scalar $\alpha$,
\begin{equation}
\rho_d^{\,\text{site}} \;=\; \alpha\cdot 1800\ \mathrm{kg\,m^{-3}},
\label{eq:alpha}
\end{equation}
with $\rho_s=1100$ kg m$^{-3}$, the e-folding depth $H$, and the published
coefficients $(A_1,A_2,B_1,B_2)$ and $k_\mathrm{am}(T)$ held fixed;
$\alpha=1$ recovers the published forward prediction. The scalar $\alpha$
is our diagnostic construction (the \citet{martinez2021} model itself is
not refit, only re-evaluated at a per-site density), and its physical
interpretation is the per-site asymptotic bulk density that
the published $K(T,\rho)$ would require to reproduce the in-situ data. The
$\alpha$ retrieval is not part of the headline $K_d$ result; it is
interpreted in \S\ref{sec:disc-prior}.

\subsection{Borestem diagnostic and the meter-scale-sensor cut}%
\label{sec:results-amp}\label{sec:borestem}
The meter-scale ($z>\zborestem$) sensor set used in the
$K_d$ retrieval is defined by a physics-based amplitude diagnostic, not by
the residuals it would produce. The amplitude of the diurnal wave serves as
a frequency-domain check independent of the lunation-mean comparison. The
per-sensor diurnal amplitude is estimated as
$\sqrt{2}\,\sigma_{\Delta t}$, where $\sigma_{\Delta t}$ is the
within-stability-window standard deviation; the proxy equals the
half-range exactly for a sinusoidal signal, is diurnal-dominated for
shallow sensors, and saturates at the digitization noise floor below the
skin depth --- precisely the contrast the diagnostic requires
(Fig.~\ref{fig:amplitude}). Two
features are diagnostic. Above $z=\zborestem$, observed amplitudes exceed
the modeled regolith-attenuation curve by an order of magnitude or more,
the signature of heat conducted along the fiberglass borestem from the
lunar surface, which thermally short-circuits the upper $\sim\zborestem$ of
regolith \citep{nagihara2018}. Below $\zborestem$, both published $K(z)$
forms predict diurnal amplitudes $\lesssim 10^{-7}$ K ($\sim$18 e-foldings
of the diurnal skin depth $\delta\sim 3$--5 cm), two to three orders of
magnitude below the 0.05--0.2 K digitization noise floor of the restored
record. Meter-scale-sensor amplitudes therefore do
not constrain the thermal diffusivity $\kappa = K/(\rho c_p)$ at the
in-situ scale; the Hayne and \citet{martinez2021} attenuation curves
plotted in Figure~\ref{fig:amplitude} agree to within $\sim$10\% in skin
depth and overlap visually below the noise floor, both indicating ``deep
$=$ unresolved''. The shallow-sensor amplitude excess, evaluated against a
prediction that depends only on $K(T,z)$, $\rho(z)$, and $c_p(T)$, defines
the meter-scale ($z>\zborestem$) validation set used in the
retrieval.

\begin{figure}[!htbp]
\centering
\includegraphics[width=\linewidth]{figures/fig_amplitude_vs_depth.pdf}
\caption{Diurnal amplitude ($\sqrt{2}\times$ the within-stability-window
standard deviation) versus sensor depth at both Apollo sites. Open circles:
borestem-zone sensors ($z<80$~cm); filled circles: meter-scale sensors;
dotted leaders mark each borestem sensor's heat-short excess above the
modeled attenuation. Solid teal: \Hayne\ attenuation; dashed violet:
\citet{martinez2021} --- visually indistinguishable on the log axis
($T=250$ K, published density 1800 kg m$^{-3}$). The order-of-magnitude
shallow-sensor excess is the borestem heat-short signature.}
\label{fig:amplitude}
\end{figure}


\subsection{Per-site $K_d$ retrieval (Hayne shape held fixed)}%
\label{sec:methods-kd}
$K_d$ is retrieved separately at each site by minimizing the
meter-scale-sensor RMSE while holding the remainder
of the Hayne functional form (eqs.~\ref{eq:K}--\ref{eq:cp}) fixed.
$(K_s,H,\chi,T_\mathrm{ref})$ and the density and $c_p$ functions are held
at their published values, and $K_d$ is swept across site-specific grids:
29 points spanning
1.0--$15.0\times10^{-3}\,\si{\watt\per\meter\per\kelvin}$ at A15 and 32
points spanning
3.0--$25.0\times10^{-3}\,\si{\watt\per\meter\per\kelvin}$ at A17, sized so
that both parabolic minima and the bootstrap upper percentiles fall inside
the swept range (Fig.~\ref{fig:bootstrap}).
 At each $K_d$ a full Crank--Nicolson spin-up is
run to convergence, and the
meter-scale-sensor RMSE is evaluated as
\begin{equation}
\mathrm{RMSE}(K_d) = \Bigl\langle\bigl[T_\mathrm{model}(z_i;K_d) - T_\mathrm{eq}^\mathrm{HFE}(z_i)\bigr]^2\Bigr\rangle_{z_i\ge\zborestem}^{1/2}.
\label{eq:rmse-kd}
\end{equation}
The point estimate $K_d^{*}$ is the parabolic minimum through the three
grid values bracketing the RMSE minimum. Its uncertainty is
quantified by a non-parametric bootstrap with $N_\mathrm{boot}=1500$
resamples of the meter-scale-sensor set drawn with
replacement, with the $\pm 2.5$ cm sensor-placement uncertainty
\citep{nagihara2018} propagated as a per-resample Gaussian depth jitter,
following standard small-sample regression practice \citep{press2007}. The
bootstrap distribution defines 95\% confidence intervals on $K_d^{*}$ and
on the inter-site contrast $\Delta K_d^{*}$, and provides the
bootstrap tail
proportion $P_{\rm boot}(\Delta K_d^{*}{\le}0)$, the fraction of
fitted (non-null-centered) resample draws with a non-positive contrast,
reported in place of a conventional null-centered $p$-value
(Fig.~\ref{fig:bootstrap}).
Percentile-method intervals are reported throughout; bias-corrected
accelerated (BCa) intervals \citep{efron1993} agree with the percentile
endpoints to within $1.3\,\si{\milli\watt\per\meter\per\kelvin}$ at both
sites and on the contrast, and do not alter any qualitative conclusion
--- a difference comparable to the CI half-widths themselves
($\sim$1.4 at A15, $\sim$1.0 at A17), so the percentile intervals
should be read as approximate at these sample sizes ($N=7$ and $16$). The
retrieval has one free parameter.

Models with different free-parameter counts on the same
meter-scale-sensor set are compared using the
small-sample-corrected Akaike information criterion \citep{burnham2002},
$\mathrm{AICc}=N\ln(\mathrm{RSS}/N)+2k+2k(k{+}1)/(N{-}k{-}1)$, where RSS is
the meter-scale-sensor residual sum of squares and
$k$ counts the fitted parameters plus one for the residual-variance
estimate ($k=1$ for fixed-$K_d$ rows, $k=2$ for each retrieval). Differences $\Delta\mathrm{AICc}$ are reported relative
to the best model at each site; following \citet{burnham2002},
$\Delta\mathrm{AICc}\lesssim 2$ indicates statistically indistinguishable
models and $\Delta\mathrm{AICc}\gtrsim 10$ a decisive preference. A reduced
$\chi^2$ is not reported: the within-stability-window scatter
($\sigma\sim 0.03$--$0.05$ K) characterizes short-term thermometer
precision rather than the model-adequacy error budget, yields
$\chi^2_\nu\gg 1$ for every model, and is therefore uninformative for model
selection in this context.

% =====================================================================
\section{Results}\label{sec:results}

\subsection{Annual-mean subsurface profile and the global-model gap}%
\label{sec:results-mean}
Both published global models
reproduce the lunation-mean borehole profiles to within $\sim$1 K, but
systematic per-site residuals remain. Figure~\ref{fig:meanT} compares
the two modeled time-averaged $\langle T(z)\rangle$ profiles,
\citet{hayne2017} at the global
$K_d=3.4\,\si{\milli\watt\per\meter\per\kelvin}$ and
\citet{martinez2021} $K(T,\rho)$ evaluated forward, against the
stability-window means. The published Hayne value runs warm relative to the
HFE record by a mean of $+0.62$ K at A15 and $+0.54$ K at A17, with RMSE 1.09 and
0.89 K --- the residual structure, not the mean offset, carries the
per-site information; the \citet{martinez2021} prediction
is
comparable (RMSE 1.05 and 0.63 K) (Table~\ref{tab:headline}). The
shaded band ($z<\zborestem$) marks the borestem-contaminated zone,
excluded \emph{a priori} from the quantitative comparison
(\S\ref{sec:borestem}). The residual structure that remains --- in particular the
under-predicted temperature gradient between the 80-cm and deepest
sensors at A17 --- motivates the per-site $K_d$ retrieval of
\S\ref{sec:results-kd}.

\begin{figure}[!htbp]
\centering
\includegraphics[width=\linewidth]{figures/fig_apollo_mean_T_profile.pdf}
\caption{Annual-mean subsurface temperature $\langle T(z)\rangle$ for the two
published models at nominal (un-retrieved) parameters --- \Hayne\ at the
global $K_d=3.4\,\si{\milli\watt\per\meter\per\kelvin}$ (solid) and
\citet{martinez2021} $K(T,\rho)$ forward (dashed) --- against the Apollo HFE
stability-window means (charcoal $=$ TG, gray $=$ TR sensors; bars =
within-window standard deviation). Shaded band: the borestem zone
($z<\SI{80}{\centi\meter}$), excluded from the RMSE.}
\label{fig:meanT}
\end{figure}

\begin{table}[t]
\centering
\small
\setlength{\tabcolsep}{5pt}
\caption{\textbf{Headline meter-scale-sensor RMSE
table.} Meter-scale-sensor (%
$z=80$--$234$ cm) validation statistics for the three forward
configurations confronted with the same Apollo HFE record: the \Hayne\ form
at its published global $K_d$, the \Hayne\ form with $K_d$ refit per site
(the retrieval of \S\ref{sec:results-kd}), and the \citet{martinez2021}
$K(T,\rho)$ model evaluated forward at each site. $K_d$ in
mW m$^{-1}$ K$^{-1}$ (95\% bootstrap CI in brackets for the Hayne site-fit
row). Bias $=$ mean(model $-$ obs). A diagnostic per-site retrieval under
the \citet{martinez2021} form is presented separately in
\S\ref{sec:disc-prior}.}
\label{tab:headline}
\begin{tabular}{l l l c c c}
\toprule
Site & Model & Fitted knob & $N$ & RMSE (K) & Bias (K) \\
\midrule
Apollo 15 & Hayne (2017), global $K_d$    & ---                              & 7  & \textbf{1.09} & $+0.62$ \\
Apollo 15 & Hayne, $K_d^{*}$ fit          & $K_d^{*}=4.60\,[4.18,6.96]$      & 7  & \textbf{1.00} & $+0.44$ \\
Apollo 15 & \citet{martinez2021} forward  & ---                              & 7  & \textbf{1.05} & $+0.54$ \\
\midrule
Apollo 17 & Hayne (2017), global $K_d$    & ---                              & 16 & \textbf{0.89} & $+0.54$ \\
Apollo 17 & Hayne, $K_d^{*}$ fit          & $K_d^{*}=7.08\,[6.16,8.07]$      & 16 & \textbf{0.40} & $+0.12$ \\
Apollo 17 & \citet{martinez2021} forward  & ---                              & 16 & \textbf{0.63} & $+0.09$ \\
\bottomrule
\end{tabular}

\end{table}

\label{sec:headline-table}

\subsection{Per-site $K_d$ retrieval}\label{sec:results-kd}
The meter-scale-sensor RMSE exhibits a well-resolved
minimum at each site:
\[
K_{d,\mathrm{A15}}^{*}=4.60^{+0.86}_{-0.21}\,[4.18,\,6.96]\,\si{\milli\watt\per\meter\per\kelvin},\qquad
K_{d,\mathrm{A17}}^{*}=7.08^{+0.53}_{-0.52}\,[6.16,\,8.07]\,\si{\milli\watt\per\meter\per\kelvin},
\]
with asymmetric $1\sigma$ bootstrap uncertainties and 95\% confidence
intervals in brackets (Fig.~\ref{fig:bootstrap}). The
meter-scale-sensor RMSE decreases from
1.09 K and 0.89 K under the published
global $K_d$ to 1.00 K and 0.40 K under the
per-site retrieval ($N=7$ at A15, $N=16$ at A17;
Table~\ref{tab:headline}); at A17 the refit halves the RMSE and removes the
gradient misfit (bias $+0.12$ K), while at A15 the improvement is
modest and the global value is not excluded by the RMSE criterion
alone --- it is excluded by the bootstrap CI, whose lower bound
(4.18) lies above 3.4.
$K_d$ is therefore the dominant identifiable lever between the Hayne form
and the in-situ HFE record once the borestem zone is excluded,
with the caveat that the per-site evidence is decisive at A17
($\Delta\mathrm{AICc}\approx 23$ against the global value) but not at
A15 alone ($\Delta\mathrm{AICc}\approx 3$ in favor of the
global value once the parameter penalty is paid; the inter-site
contrast, not the A15 refit, carries the statistical result.

The minimum is not an artifact of the sweep grid. Extending the sweep
beyond the fitting envelope reveals that each RMSE curve rises
monotonically on the high-$K_d$ side (Fig.~\ref{fig:kdsweep}), confirming
bracketed minima rather than descending shoulders. Nor is it an
artifact of the thermal-state initialization: the equilibrium
construction of \S\ref{sec:solver} is certified independent of its
initial guess to $\le 0.03$ K, and a grid-and-timestep convergence study
(halving $\Delta z_0$ and $\Delta t$) bounds the solver contribution to
$K_d^{*}$ at $\lesssim 0.05\,\si{\milli\watt\per\meter\per\kelvin}$,
an order of magnitude below the bootstrap uncertainty quoted above.

The \citet{martinez2021} model, evaluated forward at each site with its
published global coefficients, serves as the global-baseline comparison:
meter-scale-sensor RMSE of 1.05 K
at A15 and 0.63 K at A17, with mean biases of $+0.54$ K and $+0.09$ K
(Table~\ref{tab:headline}). It reproduces
both sites to within $\sim$1 K without any per-site freedom ---
comparable to the Hayne site-fit at A15, though still a factor of
$\sim$1.6 above it at A17. A diagnostic per-site retrieval under
this form is presented in \S\ref{sec:disc-prior}, where the implications
for the $K(T,\rho)$ parameterization are discussed.

\paragraph*{Sensitivity to the borestem exclusion depth.}
The retrieval is repeated with the borestem exclusion depth $z_b$ swept
over 60--90 cm. The retrieved values remain stable across all admissible
cuts: $K_{d,\mathrm{A15}}^{*}=4.6--4.9$ and
$K_{d,\mathrm{A17}}^{*}=7.2\,\si{\milli\watt\per\meter\per\kelvin}$ for
$z_b=70$, 80, and 90 cm, yielding a contrast of
2.6--3.0$\,\si{\milli\watt\per\meter\per\kelvin}$,
unchanged within rounding from the nominal value. The 60 cm cut alone
deviates
($K_{d,\mathrm{A17}}^{*}=8.8$, contrast 3.9),
because at this depth a
sensor inside the borestem-affected zone re-enters the A17 set and inflates
the retrieved gradient, precisely the contamination that the
amplitude-based cut (\S\ref{sec:borestem}) is designed to exclude. The
retrieval is therefore insensitive to $z_b$ across its physically
defensible range, and the 60 cm outlier confirms rather than undermines the
borestem diagnostic.

\begin{figure}[!htbp]
\centering
\includegraphics[width=\linewidth]{figures/fig_kd_sweep.pdf}
\caption{Per-site $K_d$ retrieval under the \Hayne\ form (all other Hayne
parameters at published values). Meter-scale-sensor RMSE($K_d$) at each site,
extended past the fitting grid so each curve rises again on the high-$K_d$
side --- the minima are genuine bracketed bowls. Stars: retrieved $K_d^{*}$
(horizontal bars: 95\% bootstrap CI); markers: forward-model evaluations.
Dashed vertical: \citet{hayne2017} global $K_d=3.4$; dotted:
\citet{feng2020} deep value $K_d=3.8$.}
\label{fig:kdsweep}
\end{figure}


Two summary statistics must be distinguished. The $K_d^{*}$ values quoted
above are parabolic point estimates of the RMSE minimum; the bootstrap
medians of the same quantities are 4.63 and 7.08$\,\si{\milli\watt\per\meter\per\kelvin}$
(Table~\ref{tab:errorbudget}). The inter-site contrast is
defined as the bootstrap median of the paired within-resample difference,
$\Delta K_d^{*} = 2.31^{+0.68}_{-0.90}\,\si{\milli\watt\per\meter\per\kelvin}$
($1\sigma$, asymmetric; 95\% CI %
$[-0.12,\,3.56]$, $P_{\rm boot}(\Delta K_d^{*}{\le}0)\approx 0.031$;
Fig.~\ref{fig:bootstrap}). The two sites
share no sensors, so the contrast draws combine independent
per-site resamples; an all-pairs convolution of the two bootstrap
distributions reproduces the paired interval to rounding precision. The
median-of-differences construction is also the reason
2.31 differs slightly from the difference of the point
estimates ($7.08-4.60=2.48$). The
A15 upper tail reflects a weakly constrained
shoulder in RMSE$(K_d)$ at A15 for $N=7$ meter-scale %
sensors, even though the central parabolic minimum is firmly at
4.60. The
reported $1\sigma$ intervals propagate sensor-set composition, the
$\pm 2.5$ cm depth-placement uncertainty \citep{nagihara2018}, and the
parabolic curve-fit error; component-by-component impacts of the remaining
Hayne inputs are tabulated in Table~\ref{tab:errorbudget}. Both retrievals
reject the global $K_d=3.4\,\si{\milli\watt\per\meter\per\kelvin}$ of
\citet{hayne2017} at the 95\% level (bootstrap CI lower bounds 4.18
and 6.16).

\begin{figure}[!htbp]
\centering
\includegraphics[width=\linewidth]{figures/fig_bootstrap.pdf}
\caption{\textbf{Bootstrap distributions} ($N_\mathrm{boot}=1500$,
meter-scale-sensor resampling with $\pm 2.5$\,cm depth jitter).
(a)~Per-site $K_d^{*}$ histograms (95\% CIs in the legend); dashed vertical:
the \citet{hayne2017} global value. (b)~Inter-site contrast
$\Delta K_d^{*}\!=\!K_d^{\rm A17}\!-\!K_d^{\rm A15}$; 95\% CI shaded, lower
bound just below zero ($p\!\approx\!0.031$).}
\label{fig:bootstrap}
\end{figure}

Applied at each site, the retrieved $K_d^{*}$ values yield
meter-scale-sensor profiles that track the HFE
observations within the $\sim 1$ K spread of the data
(Fig.~\ref{fig:thermal-profiles}). The \citet{martinez2021} forward,
plotted alongside, matches the Hayne site-fit RMSE at A15 to within
$\sim$5\% but remains a factor of $\sim$1.6 above it at A17.

\begin{figure}[!htbp]
\centering
\includegraphics[width=\linewidth]{figures/fig_thermal_profiles.pdf}
\caption{Annual-mean $T(z)$ at the two sites under the two models, against the
Apollo HFE meter-scale sensors. Solid teal: \Hayne\ at the per-site retrieved
$K_d^{*}$; dashed violet: \citet{martinez2021} forward (no fitted knob). Top:
full profile; bottom: meter-scale-sensor zoom. Colored markers: sensors used
in retrieval; gray: borestem-affected shallow sensors (excluded).}
\label{fig:thermal-profiles}
\end{figure}

A full per-component error budget for $K_d^{*}$ is deferred to
\S\ref{sec:disc-degeneracy} (Table~\ref{tab:errorbudget}), where the
basal-flux degeneracy that dominates the budget is discussed.

\subsection{Consistency checks}\label{sec:consistency}
Three diagnostics test the per-site $K_d^{*}$ values. Their status as
independent or internal is identified explicitly.

An independent Diviner surface-temperature closure test
(Fig.~\ref{fig:diviner}, discussed in \S\ref{sec:disc-prior}) shows that
both conductivity models reproduce the canonical lunar diurnal shape,
confirming that the diurnal forcing chain is correctly captured. Two
further internal checks follow.

\emph{(i) TG vs.\ TR cross-prediction.} Fitting $K_d$ on the
gradient-bridge sensors (TG) alone and predicting the ring-bridge (TR)
sensors, and vice versa, yields held-out RMSEs within
$\sim$0.11 K of the
in-sample values at both sites.
\emph{(ii) Leave-one-deepest-out.} Withholding the deepest sensor
($z=139$ cm at A15, $z=234$ cm at A17) and refitting $K_d$ on the remainder
predicts the held-out temperature to within
$|\Delta|=0.37$ K and
$|\Delta|=0.25$ K respectively. Both internal checks
confirm that the
retrieval is not driven by a single sensor or sensor type; neither
constitutes an independent test of the inter-site contrast.

\emph{Bayesian cross-check.} The non-parametric bootstrap is reported as
the primary uncertainty estimate throughout because it makes no assumption
about the basal heat flux and propagates only quantities measured at the
boreholes (sensor-set composition and placement depth). As an independent
cross-check on the \emph{direction} of the contrast, not as a second
estimate of its magnitude, we also ran a Bayesian inversion propagating the
\citet{saito2007,nagihara2018} prior on $Q_b$
(Fig.~\ref{fig:robustness}a). Its marginal MCMC medians
are $K_{d,\mathrm{A15}}=4.3$ and
$K_{d,\mathrm{A17}}=8.3$
mW m$^{-1}$ K$^{-1}$ (contrast median $+3.9$).
The corresponding $Q_b$ marginals pull below their priors at both
sites --- medians 17.8 (A15, prior $21\pm3.2$) and 13.9 mW m$^{-2}$
(A17, prior $16\pm2.4$) --- i.e., the HFE record itself mildly favors
the lower, \citet{saito2007}-direction fluxes, independently at the two
boreholes; a firmer per-site $Q_b$ remains the single most valuable
external input this retrieval could receive. These
differ from the
bootstrap point estimates (4.60,
7.08; contrast $+2.5$) by construction,
not by disagreement: the $Q_b$ prior slides each marginal along the
$K_d/Q_b$ degeneracy ridge (\S\ref{sec:disc-kd}), so the two methods are
not independent measurements of the absolute $K_d$. The ordering not only survives
the $Q_b$ smearing --- it is insensitive to it: on the directly solved
RMSE$(K_d,Q_b)$ surface (no proportional-degeneracy surrogate;
\S\ref{sec:disc-degeneracy}), $P(K_d^{\rm A17}>K_d^{\rm A15})$ stays at
$99.2$--$99.4\%$ for Gaussian $Q_b$ priors from $\pm$10\% to
$\pm$40\% (archived scan: \texttt{results/qb\_prior\_width\_scan.json}),
and the contrast median grows with prior width (from
$+3.4$ to $+4.8$ mW m$^{-1}$ K$^{-1}$) because the two sites respond to a
flux revision in opposite directions. The Bayesian and bootstrap views
therefore bracket the honest statement: conditional on the published
basal fluxes the contrast magnitude is only marginally resolved
($P_{\rm boot}{\approx}0.03$, bootstrap 95\% interval spanning zero), while once the
basal-flux debate itself is propagated the \emph{ordering} becomes
decisive (99\%, with the credible interval on the contrast excluding
zero) --- though its absolute magnitude then inherits the full flux
uncertainty (\S\ref{sec:disc-degeneracy}). Both statements are
conditional on the assumed per-sensor noise scale (0.5 K, chosen at the
scale of the residual model--data RMSE of Table~\ref{tab:headline}
rather than the 0.03--0.05 K within-window thermometer precision, which
characterizes short-term stability only) that sets the
likelihood sharpness.%
% Q_b prior-width scan, Jacobian-corrected posterior (2026-07-13, codex F7
% fix; reproducible artifact: pipeline/compute/qb_prior_width_scan.py ->
% results/qb_prior_width_scan.json): P(A17>A15) = 99.21/99.25/99.34/99.42 %
% at prior sigma = 10/15/25/40 % of Q_b; contrast median +3.43/+3.92/+4.49/
% +4.79, 95% CrI low edge +0.63/+0.75/+0.93/+1.06.
\paragraph*{Common-epoch sensitivity.}
The stability windows end at different calendar epochs while the
record still warms: at A17 the retained sensors at 130--178 cm carry
mid-1977 windows while every sensor below 180 cm carries a 1974 window
(a 3.1-yr spread aligned with depth); at A15 the spread is 1.0 yr with
trailing slopes $\le$0.22 K yr$^{-1}$. Rebuilding every retained
sensor's $T_\mathrm{eq}$ from one shared 1974 calendar window (every
sensor has data there --- no extrapolation) leaves A15 unchanged
($|\Delta K_d^{*}|\le 0.01$) but moves A17 to
$K_d^{*}=6.16$ mW m$^{-1}$ K$^{-1}$ ($-0.89$; bootstrap 95\% CI
$[5.92, 6.81]$) --- and improves the fit (RMSE
$0.40\rightarrow0.36$ K), as expected if the mixed-epoch profile was
internally inconsistent. A slope-translated construction agrees to
0.03 mW m$^{-1}$ K$^{-1}$. The contrast under the common epoch is
$+1.55$ mW m$^{-1}$ K$^{-1}$ (95\% CI $[-0.11, 2.31]$,
$P_{\rm boot}(\Delta K_d^{*}{\le}0)=0.04$): still positive, still
rejecting the global value at both sites, but more marginal than at the
certified windows. Epoch mixing is therefore the largest single
methodological systematic at A17 and is carried as $\sigma_{\rm epoch}$
in Table~\ref{tab:errorbudget}
(archived: \texttt{results/common\_epoch\_sensitivity.json}).

The bootstrap interval
on the contrast
magnitude ($\Delta
K_d^{*}=2.3$, 95\% CI $[-0.1,3.6]$) is carried forward
as the headline measurement.


% =====================================================================
\FloatBarrier
\section{Discussion}\label{sec:discussion}

\subsection{The two boreholes require different meter-scale temperature gradients}%
\label{sec:disc-kd}
The HFE record does not admit a single meter-scale %
conductivity that fits both sites. Under the Hayne form at the published
basal heat fluxes, the value that fits the A17
meter-scale-sensor residual
(7.08) degrades
A15 to a $+0.70$ K bias (RMSE 1.16 K), and the value that
fits A15 (4.60) degrades the A17 RMSE from 0.40 to 0.54 K. The
published global value of \citet{hayne2017}, applied
unchanged, runs warm at both sites ($+0.62$ K at A15, $+0.54$ K at A17;
Table~\ref{tab:headline}) and, more importantly, under-predicts the
A17 meter-scale temperature gradient. The
per-site retrievals are $K_d^{\mathrm{A15}}=4.60$ and
$K_d^{\mathrm{A17}}=7.08\,\si{\milli\watt\per\meter\per\kelvin}$ (95\%
bootstrap CIs $[4.18,\,6.96]$
and $[6.16,\,8.07]$). The lower 95\% bound
on the contrast distribution (Fig.~\ref{fig:bootstrap}b)
spans zero
($\Delta K_d^{*}=2.31$, 95\% CI $[-0.12,\,3.56]$,
$P_{\rm boot}(\Delta K_d^{*}{\le}0)\approx 0.031$):
conditional on the published basal heat fluxes, the \emph{requirement} of
different meter-scale temperature gradients is established (decisively at
A17, $\Delta\mathrm{AICc}\approx 23$), while the \emph{magnitude} of the
conductivity contrast is favored by $\sim$97\% of draws but not resolved
at the 95\% level. The conditionality matters --- under the $K_d/Q_b$
degeneracy of \S\ref{sec:disc-degeneracy} the attribution to
conductivity rather than basal flux still rests on the $Q_b$
assumption.

The pooled three-model comparison (Table~\ref{tab:modelcomp})
has, at the certified
numerics, no discriminating power at all: all three models --- per-site
$K_d$ (M1), shared $K_d$ with free per-site $Q_b$ (M2), and the strict
shared-$K_d$ null (M3) --- fall within $\Delta\mathrm{AICc}\le 0.5$ of
one another, a statistical tie. This is expected rather than
troubling: the pooled statistic dilutes the A17 signal with the
less-informative A15 set ($N=7$), whereas the per-site comparison is
decisive at A17 ($\Delta\mathrm{AICc}\approx 23$ for the site fit over
the global value) and the bootstrap favors a positive
contrast ($P_{\rm boot}{\approx}0.03$; its 95\% interval spans zero). The
data therefore favor a per-site difference in
meter-scale temperature gradient without excluding a
uniform-conductivity alternative in which the two basal heat fluxes differ.

\begin{table}[!htbp]
\centering
\small
\caption{Pooled meter-scale-sensor model comparison
($N=23$: $7$ at A15, $16$ at A17). M1 retrieves a per-site $K_d$; M2 shares
one $K_d$ and frees the per-site $Q_b$ within the published envelope; M3
fixes both. The three models are compared on a shared coarse
($1$\,mW) $K_d$ grid, so their tabulated $K_d$ are rounded relative to
the dense per-site headline retrieval ($4.60/7.08$); the pooled test
probes model \emph{structure}, not the vertex location. $\Delta$AICc
is relative to the best model. $K_d$ in
mW\,m$^{-1}$\,K$^{-1}$, $Q_b$ in mW\,m$^{-2}$.}
\label{tab:modelcomp}
\begin{tabular}{llccc}
\toprule
Model & $K_d$ (A15\,/\,A17) & $Q_b$ (A15\,/\,A17) & RMSE (K) & $\Delta$AICc \\
\midrule
M1\ \ variable $K_d$, $Q_b$ fixed & $5.0$\,/\,$7.0$ & $21$\,/\,$16$ (fixed) & 0.64 & $+0.1$ \\
M2\ \ uniform $K_d$, $Q_b$ free   & $7.0$ (shared)   & $14$\,/\,$15$ (fitted) & 0.61 & $+0.5$ \\
M3\ \ uniform $K_d$, $Q_b$ fixed  & $6.0$ (shared)  & $21$\,/\,$16$ (fixed) & 0.68 & \textbf{0.0} \\
\bottomrule
\end{tabular}
\end{table}

\subsection{The retrieved absolute $K_d$ values are degeneracy-limited}%
\label{sec:disc-degeneracy}
The absolute $K_d^{*}$ values must be interpreted as conditional on the
adopted basal heat flux, not as free-standing measurements. At leading order the meter-scale
temperature gradient $|\partial_z T|=Q_b/K_d$ suggests the familiar
proportional degeneracy $(K_d,Q_b)\!\to\!(\alpha K_d,\alpha Q_b)$; the
retrieval, however, also matches the \emph{absolute} sensor temperatures,
whose offset through the low-conductivity skin column depends on $Q_b/K$
non-proportionally --- a direct test (rescaling $Q_b$ and $K_d$ together
by $10/16$ at A17 shifts the whole sensor band by a depth-uniform
$-1.3$~K) shows the invariance fails, so the realized trade-off must be
mapped by re-retrieval rather than assumed. The degeneracy is mapped
quantitatively in Fig.~\ref{fig:robustness}a and is visible in the joint
Bayesian posteriors. Adopting the canonical
$Q_b^{\mathrm{A17}}=16\,\si{\milli\watt\per\meter\squared}$
\citep{langseth1976,nagihara2018} yields $K_d=7.08$;
re-retrieving at the
lower $\sim$10 mW m$^{-2}$ value of \citet{saito2007,nagihara2018} moves
the minimum up to $\sim$10.1 mW m$^{-1}$ K$^{-1}$ --- and deepens
it (RMSE 0.29 vs 0.40 K) --- while the upper bound $\sim$18 mW m$^{-2}$
lowers it to $\sim$6.8 mW m$^{-1}$ K$^{-1}$.
The A17 $K_d^{*}$ therefore carries a
$\sim$$^{+43}_{-4}\%$ envelope determined
entirely by the basal-flux debate. A15
responds in the opposite sense --- its steeper gradient dominates
the fit, so $K_d^{*}$ falls from 5.4 to 2.7 as $Q_b^{\mathrm{A15}}$ is
revised from $+25\%$ to $-25\%$ of the published 21 mW m$^{-2}$. Because
the two sites respond oppositely, the inter-site \emph{contrast} is more
robust to the flux debate than either absolute value: it remains positive
across the entire debated envelope even with the two fluxes varied
independently (worst pairing, both fluxes at their upper bounds: $6.8-5.3=+1.5$ mW m$^{-1}$ K$^{-1}$). What
the boreholes cannot deliver is each site's \emph{absolute} $K_d$ --- and
fit quality alone cannot arbitrate the flux debate, since at A17 the
\citet{saito2007} flux actually fits the HFE record slightly better than
the published value.
This represents the intrinsic limit of what two boreholes can deliver.

\begin{table}[!htbp]
\centering
\footnotesize
\setlength{\tabcolsep}{3.4pt}
\renewcommand{\arraystretch}{1.10}
\caption{\textbf{Per-component sensitivity budget for $K_d^{*}$.}
Entries are half-range sensitivity indices from scenario sweeps, not
calibrated Gaussian standard deviations; the quadrature total is an
independence-idealized summary (see text).
Values in mW m$^{-1}$ K$^{-1}$; entries below 0.05 are reported as
$<\!0.05$. All components are re-derived with the certified
equilibrium solver. The quadrature total covers the statistical and
measured-input components; $\chi$ and $H$ parameterize the
\citet{hayne2017} functional form that the retrieval holds fixed, and
are reported separately as conditionality terms (see text) rather than
summed --- re-retrieving at the \citet{vasavada2012} normalization
$\chi=1.48$ drives the RMSE minimum outside the physical sweep range at
both sites.}
\label{tab:errorbudget}
\begin{tabular}{@{}l p{0.19\linewidth} p{0.27\linewidth} r r@{}}
\toprule
Component & Envelope & Method & A15 & A17 \\
\midrule
\multicolumn{5}{@{}l}{\textit{Statistical and solver}} \\[2pt]
$\sigma_\mathrm{stat}$
  & $\pm 2.5$\,cm depth jitter
  & bootstrap 16--84\,\% half-range
  & 0.53 & 0.52 \\
$\sigma_\mathrm{solver}$
  & grid \& $\Delta t$ halved
  & convergence study, re-retrieve $K_d^{*}$ (\S\ref{sec:solver})
  & $<\!0.05$ & $<\!0.05$ \\
\midrule
\multicolumn{5}{@{}l}{\textit{Measured inputs}} \\[2pt]
$\sigma_{Q_b}$
  & $Q_b$\,A15$\in\![14,25]$,\,A17$\in\![10,18]$
  & direct $K_d^{*}(Q_b)$ re-retrievals, half-range (\S\ref{sec:disc-degeneracy}; asymmetric: $^{+0.7}_{-2.5}$\,/\,$^{+3.1}_{-0.3}$)
  & 1.61 & 1.66 \\
$\sigma_A$\,(albedo)
  & $A\!\pm\!0.01$
  & re-retrieve $K_d^{*}$ at perturbed $A$
  & 0.72 & 2.96 \\
$\sigma_{K_s}$
  & $K_s\!\pm\!30\%$
  & re-retrieve at perturbed $K_s$
  & 0.36 & 1.55 \\
$\sigma_\rho$
  & $\rho_d\!\in\![1700,2000]$\,kg\,m$^{-3}$
  & re-retrieve at perturbed $\rho_d$ ($\rho$ enters $\rho c_p$ only)
  & $<\!0.05$ & 0.33 \\
\midrule
\multicolumn{5}{@{}l}{\textit{Methodological choices (this work)}} \\[2pt]
$\sigma_\mathrm{thr}$
  & $0.04$--$0.16$\,K/yr ($4\!\times$ adopted)
  & threshold sweep, half-range of $K_d^{*}$
  & $<\!0.05$ & $<\!0.05$ \\
$\sigma_{z_b}$
  & $z_b\!\in\!\{70,80,90\}$\,cm
  & cut sweep, half-range of $K_d^{*}$
  & 0.12 & $<\!0.05$ \\
$\sigma_\mathrm{epoch}$
  & common-1974 vs certified windows
  & re-fit at common-epoch $T_\mathrm{eq}$ (\S\ref{sec:consistency}; one-sided: A17 $-0.89$)
  & $<\!0.05$ & 0.89 \\
\midrule
\textbf{Total (quadrature)} & & & \textbf{1.88} & \textbf{3.88} \\
Bootstrap median $K_d^{*}$  & & & 4.63 & 7.08 \\
\midrule
\multicolumn{5}{@{}l}{\textit{Conditionality of the held-fixed form (reported, not summed)}} \\[2pt]
$\chi$
  & $\chi\!\in\!\{1.48, 2.7\}$
  & re-retrieve at the \citet{vasavada2012} endpoint
  & \multicolumn{2}{c}{minimum exits sweep ($\le$1.0\,/\,$\ge$25)} \\
$H$
  & $H\!\in\![3,10]$\,cm
  & per-$H$ $K_d^{*}$ half-range along the valley
  & 0.51 & 3.47 \\
\bottomrule
\end{tabular}
\end{table}

\paragraph*{Reading the per-component error budget.}
The derivation of each row in Table~\ref{tab:errorbudget} falls into one
of three categories: $\sigma_\mathrm{stat}$ is the 16--84 percentile
half-width of the bootstrap distribution under the $\pm 2.5$ cm
astronaut-emplacement envelope reported by \citet{nagihara2018};
$\sigma_{Q_b}$ is obtained in closed form through the $K_d/Q_b$ degeneracy
across the published reduction range (\citealp{langseth1976} canonical;
\citealp{saito2007,nagihara2018} reanalysis); every remaining row is
computed by re-running the full retrieval at the perturbed input, with the
envelopes taken from the published literature ($\sigma_A$: the
\citet{vasavada2012} lat-band albedo scatter of $\pm0.01$;
$\sigma_{K_s}$: the \citet{cremers1971} laboratory scatter on Apollo~11
fines; $\sigma_\rho$: the \citet{mitchell1973,carrier1991} Apollo-core
compaction range; $\sigma_\mathrm{solver}$: the certified residual of the
equilibrium construction, \S\ref{sec:solver}).

\emph{How the total is combined.} The \textbf{Total (quadrature)} row is
the independent-errors combination $\sigma_{\rm tot}=\sqrt{\sum_i
\sigma_i^2}$ over the statistical, solver, measured-input, and
methodological rows. The worked arithmetic is
\begin{equation*}
\sigma_{\rm tot}^{\rm A15} =
\big(0.53^{2}+0.04^{2}+1.61^{2}+0.72^{2}+0.36^{2}+0.01^{2}+0.01^{2}+0.12^{2}\big)^{1/2}
= 1.88,
\end{equation*}
\begin{equation*}
\sigma_{\rm tot}^{\rm A17} =
\big(0.52^{2}+0.04^{2}+1.66^{2}+2.96^{2}+1.55^{2}+0.33^{2}+0.03^{2}+0^{2}+0.89^{2}\big)^{1/2}
\approx 3.88~\text{mW m}^{-1}\,\text{K}^{-1}.
\end{equation*}
We flag one residual correlation: $\sigma_\mathrm{stat}$ and $\sigma_{Q_b}$
both propagate along the $K_d/Q_b$ degeneracy ridge of
Fig.~\ref{fig:robustness}a, so the totals are independence-idealized
rather than fully decorrelated standard errors; more generally,
several rows are half-ranges of strongly asymmetric or nonlinear
responses (the A17 albedo row especially), so the totals are sensitivity
summaries, not calibrated $1\sigma$ uncertainties.

\emph{The conditionality terms.} The $\chi$ and $H$ rows are deliberately
excluded from the quadrature: they are not measured inputs with Gaussian
envelopes but coefficients of the \citet{hayne2017} functional form that
the retrieval holds fixed. Re-retrieving at the \citet{vasavada2012}
normalization $\chi=1.48$ drives the RMSE minimum outside the physical
sweep range at both sites --- under the converged thermal physics, the
rectified lunation-mean temperature at sensor depths depends strongly on the
radiative coefficient, so the Hayne form at the Vasavada $\chi$ cannot
reproduce the HFE record at any admissible $K_d$. $K_d^{*}$ is therefore
reported conditional on the published $(\chi, H)=(2.7, 6\,\mathrm{cm})$,
in precisely the same sense as it is conditional on the published $Q_b$;
the inter-site \emph{contrast} is far less exposed, since both sites share
the same form.

Three conclusions follow. First, both budgets are now led by a single
row: at A15 the direct-map $\sigma_{Q_b}$ alone carries $\sim$73\% of
the variance (the flux debate, not the sensors, is the accuracy limit),
whereas the A17 budget is dominated by
$\sigma_A$ and $\sigma_{K_s}$: at A17 the lunation-mean response to $K_d$ is
nearly flat over part of the sweep range, so even the $\pm0.01$ albedo
envelope moves the minimum substantially --- the statistical precision of
the A17 retrieval (CI half-width $\sim$1.0) is real, but its accuracy is
systematics-limited at the $\pm$3.8 level (the quadrature total less its
statistical row). Second, $\sigma_\rho$ is
structurally suppressed under the Hayne $K(T,z)$ form (eq.~\ref{eq:K}),
which contains no explicit $\rho$ dependence: $\rho(z)$ enters only the
volumetric heat capacity $\rho c_p$ and perturbs the rectified lunation-mean
column at the $\lesssim$0.4 mW m$^{-1}$ K$^{-1}$ level. Third, the
methodological-choice rows $\sigma_\mathrm{thr}$ and $\sigma_{z_b}$ sit
one to two orders of magnitude below the data-side terms, indicating
robustness against both discretionary cuts; the exception is
$\sigma_\mathrm{epoch}$ at A17, the epoch-mixing systematic quantified in
\S\ref{sec:consistency}, which rivals the statistical and $K_s$ rows.

\paragraph*{Behavior of the inter-site contrast under independent $Q_b$ revision.}
How the contrast responds to a basal-flux revision cannot be read off a
proportional-rescaling rule --- the $(\alpha K_d,\alpha Q_b)$ invariance
fails for this retrieval (\S\ref{sec:disc-degeneracy}) --- so we map it
explicitly by re-retrieving $K_d^{*}$ at each site across the full
adopted reanalysis envelope ($Q_{b,\mathrm{A15}}\in[14,25]$,
$Q_{b,\mathrm{A17}}\in[10,18]$ mW m$^{-2}$; Fig.~\ref{fig:robustness}a,
one direct re-retrieval per grid point).
Because the two sites respond in opposite directions
(A15's $K_d^{*}$ rises with its flux, A17's falls), the corners of the
envelope reorganize relative to the proportional picture. The
favorable corner is both fluxes at their lower bounds
($Q_{b,\mathrm{A15}}=14$, $Q_{b,\mathrm{A17}}=10$ mW m$^{-2}$), where
the contrast expands to $+8.1$ mW m$^{-1}$ K$^{-1}$
($\sim$8.6$\sigma$); the most adverse corner is both at their upper
bounds ($25$, $18$), where it compresses to $+1.5$
($\sim$1.6$\sigma$) --- reduced, but not erased. A $-33\%$ revision of
$Q_{b,\mathrm{A17}}$ alone raises the contrast to $\sim$$+5.2$.
The conductivity contrast is therefore only marginally resolved at the
published basal fluxes ($p\approx 0.03$, its 95\% interval spanning
zero), but its sign survives every admissible independent
basal-flux assignment in the published envelope: A17 remains the more
conductive site at every grid point of the direct map.

\paragraph*{Contrast under the remaining leading systematics.}
The same direct-map logic extends to the two systematics that
rival $\sigma_{Q_b}$ in Table~\ref{tab:errorbudget}. The sites again
respond in opposite directions: raising the Bond albedo by 0.01 lowers
$K_d^{*}$ at A15 (4.60$\rightarrow$3.20) but raises it at A17
(7.08$\rightarrow$11.81), so the contrast remains positive across the
full $\pm$0.01 envelope even with the two albedos perturbed
independently (most adverse pairing $+1.3$ mW m$^{-1}$ K$^{-1}$); the
$\pm$30\% $K_s$ envelope behaves identically (most adverse pairing
$+1.6$). Like the basal-flux debate, these systematics limit each
site's \emph{absolute} $K_d^{*}$ without threatening the sign of the
contrast.

A joint $(K_d, H)$ fit over the swept range $H
= 3$--$10$ cm
(Fig.~\ref{fig:robustness}b,c) confirms that $H$ is effectively
unconstrained: $K_d$
and $H$ trade off along a shallow valley: at A15 the interior minimum,
$(4.3\,\mathrm{mW\,m^{-1}\,K^{-1}},\,8\,\mathrm{cm})$, improves on the
1-parameter retrieval at the published $H=6$ cm by only 0.10 K of RMSE;
at A17 the interior minimum,
$(11.3\,\mathrm{mW\,m^{-1}\,K^{-1}},\,8\,\mathrm{cm})$, improves on the
$H=6$ slice by only 0.09 K. The reported $K_d^{*}$ is
therefore conditional on the published $H$, and the $K_d$--$H$
trade-off is propagated into the error budget as $\sigma_H$
(Table~\ref{tab:errorbudget}). $K_d$, not
$H$, is the identified
parameter at fixed $H$; the pair is jointly identified only up to
the valley direction.

\begin{figure}[!htbp]
\centering
\includegraphics[width=\linewidth]{figures/fig_robustness.pdf}
\caption{Robustness of the per-site $K_d^{*}$ contrast. \textbf{(a)}
Contrast $\Delta K_d^{*}$ across the two basal fluxes, one direct
re-retrieval per grid point (no proportional-degeneracy assumption;
\S\ref{sec:disc-degeneracy}); key features are labelled on the map and the
dashed contours are bootstrap-$\sigma$ significance. The MCMC posterior
pulls both fluxes below their published values, yet the contrast stays
positive throughout the adopted $Q_b$ envelope (minimum $+1.5$\,mW m$^{-1}$
K$^{-1}$; it dips negative only in the far, implausible
low-$Q_{b,\mathrm{A15}}$ corner), and the ordering
$K_d^{*}(\mathrm{A17}) > K_d^{*}(\mathrm{A15})$ holds at $\approx$99.2\%.
\textbf{(b,c)} Joint $(K_d, H)$ RMSE surface per site; the A17 sweep, now
extended to 12 mW m$^{-1}$ K$^{-1}$, reaches an interior valley bottom near
$11.3$ mW m$^{-1}$ K$^{-1}$ at $H\approx8$ cm (coral diamond). Teal circle:
1-parameter retrieval at $H=6$ cm.}
\label{fig:robustness}
\end{figure}

\subsection{Comparison with Mart{\'\i}nez \& Siegler (2021) and prior estimates}%
\label{sec:disc-prior}
The second globally calibrated
form fares notably better. The \citet{martinez2021} $K(T,\rho)$ model,
with its published coefficients and no per-site freedom, yields a
meter-scale-sensor RMSE of 0.99 K at A15 and 0.64 K at A17
(Table~\ref{tab:headline}, Fig.~\ref{fig:thermal-profiles}) ---
comparable to the Hayne site-fit at A15 and within a factor of
$\sim$1.9 of it at A17. Its explicitly density- and
temperature-dependent decomposition evidently transfers from the
diurnal-skin calibration to the meter-scale regolith more faithfully
than the Hayne smooth-exponential at its global $K_d$. The per-site
Hayne retrieval nonetheless remains the sharper instrument at A17,
where it halves the residual RMSE.

\paragraph*{Diagnostic: Martinez per-site density retrieval.}
The fairness of the forward Martinez comparison can be tested by granting
the published $K(T,\rho)$ one per-site lever and asking whether that lever
can absorb the A17 residual within physical bounds. We note at the outset
that \citet{martinez2021} calibrated their four coefficients against surface
brightness-temperature observations and did not validate the extrapolation
to depths of order one meter; the test we run here is therefore a check of
whether the published surface-calibrated form, when extended to the
meter-scale regolith, can fit the in-situ HFE data with one
per-site lever (the asymptotic bulk density entering the
model). We do so under the construction of \S\ref{sec:methods-kd}
(eq.~\ref{eq:alpha}): rescale the asymptote of the Hayne-form
$\rho(z)$ entering equation~(\ref{eq:K-ms}) by a single per-site scalar
$\alpha$, leaving the four Martinez coefficients fixed, and sweep $\alpha$
against the same meter-scale-sensor RMSE. The
retrieved values are $\alpha^{*}_{\mathrm{A15}}=1.19$
(asymptotic bulk density
$\rho_d^{*}\approx 2139$
kg m$^{-3}$, RMSE 0.91 K) and
$\alpha^{*}_{\mathrm{A17}}=1.04$
($\rho_d^{*}\approx 1877$ kg m$^{-3}$,
RMSE 0.60 K); at A15 the density
lever slightly betters the Hayne site-fit RMSE (0.91 vs 1.00 K), while
at A17 the bowl bottoms near the \citet{hayne2017} density and the
lever buys little ($0.63\rightarrow0.60$ K). The physical interpretation of
$\alpha$ is the asymptotic bulk density that the published
Martinez $K(T,\rho)$ would require to reproduce each site's in-situ profile.

The A15 retrieval sits $\sim$7\% above the Apollo-core compaction ceiling
of $\sim$2000 kg m$^{-3}$ \citep{mitchell1973,carrier1991} --- outside the
core envelope, though far below any solid-rock bound, and best read as
marginal rather than comfortable. The A17 retrieval bottoms
at $\rho_d^{*}\approx 1877$ kg m$^{-3}$, comfortably inside the
Apollo-core envelope (1700--2000 kg m$^{-3}$): the published
$K(T,\rho)$ is physically admissible at both boreholes.
Figure~S1 (Supporting Information) visualizes the result --- both
RMSE($\alpha$) bowls bottom inside or at the edge of the admissible
band. The flat A17 bowl also exposes the limitation: because the bowl
bottoms near the published density, no admissible density adjustment
closes the remaining factor-of-$\sim$1.5 RMSE gap to the per-site
Hayne fit; sharpening the A17 gradient further requires freedom in the
conductivity mechanism itself (the contact/radiative split), not in
the density input. A full per-site recalibration of the
\citet{martinez2021} coefficients against the HFE record is the
natural follow-on study.

% fig_alpha_sweep moved to Supporting Information (Figure S1) in the
% 2026-07-14 figure-cosmetics pass: it is a diagnostic sensitivity check,
% not a headline result, and demoting it brings the main text closer to
% the AGU publication-unit guideline. Main-text references now point to
% "Figure S1 (Supporting Information)".

The inferred ratio
$K_d^{\mathrm{A17}}/K_d^{\mathrm{A15}}\approx 1.5$ is
physically plausible if interpreted as a conductivity
effect. Two
mechanisms act in the appropriate direction and jointly cover the
required factor. First, Apollo cores span bulk
porosities $\phi\sim 0.30$--0.45 \citep{mitchell1973,carrier1991}; the
effective-medium scaling $K_\mathrm{eff}\propto(1-\phi)^2$ of
\citet{wood2020} accounts for a factor of $\sim$1.5 change at the upper end
of that range. Second, solid lunar basalt is $\sim$30\% more conductive
than solid highland anorthosite \citep{horai1972}, and A17 lies in a
basaltic embayment while A15 occupies a mixed mare/highland margin
\citep{korotev2005,wieczorek2006,crawford2015}; the product of the
two mechanisms ($\sim$1.9) comfortably covers the retrieved ratio
without requiring either at its extreme.

This is the physical sense in which the result is realistic despite
the steady-state model/data distinction. The solver represents the
asymptotic regolith column that would exist under repeated solar forcing
and the basal heat flux; the Apollo late tails are imperfect observations
of that same asymptote because the probes were still recovering from
emplacement. Residual drift can bias individual $T_\mathrm{eq}$ estimates
by fractions of a kelvin, so it belongs in the systematic error budget, but
it does not require an unphysical conductivity: the retrieved
$K_d$ values remain in the milliwatt-per-meter-per-kelvin range expected
for porous lunar regolith, and the A17/A15 ratio is covered by plausible
porosity and lithology differences.

Prior in-situ estimates on the same HFE record provide the sharpest
external reference, provided the comparison is made at matching
definition. The original in-situ reductions reported
\emph{total} effective conductivities of the compacted regolith below
$\sim$10 cm of $10\pm1$ (A15) and $15\pm1.5$ (A17)
mW m$^{-1}$ K$^{-1}$ \citep{langseth1976}.
% provenance: ALSEP Final Technical Report summary table (p.5): 'Average
% conductivity below 10 cm, mW/m-K: 10 +/-10% (A15), 15 +/-10% (A17)';
% verified from code/references/lunar_hfe_study_pack, audit 2026-07-13.
Compared at that common definition, the present retrieval
reproduces the historical in-situ values: the effective
$K(T,z)$ at the retained sensors, evaluated at $K_d^{*}$, is 9.3 (A15)
and 14.6 (A17) mW m$^{-1}$ K$^{-1}$ --- within $\sim$7\% of
\citet{langseth1976} at both sites, including the A17 $>$ A15 ordering.
The contribution of this work is therefore not an upward revision of the
in-situ conductivity but its re-parameterization: a Diviner-consistent
$K(T,z)$ form with a single per-site free parameter, a defensible
contact/radiative decomposition, and a quantified uncertainty budget;
subsequent reanalyses of the reduction assumptions
\citep{keihm1984,grott2010} are subsumed rather than extended.

\paragraph*{Surface-temperature consistency check.}
The Diviner GCP is external to the retrieval (it enters neither the
bootstrap nor the meter-scale-sensor RMSE), so it
provides an independent check on the diurnal forcing chain
(Fig.~\ref{fig:diviner}). Both parameterizations reproduce the canonical
lunar diurnal shape, with full-cycle RMSE of 4.1\,K at A15 and 6.3\,K at A17 for the Hayne
site fit, and 3.6\,K at both sites for the Martinez forward. Both
biases are dominated by the anisothermal-footprint and
sub-pixel rock-population effects inherent to Diviner brightness
temperatures \citep{bandfield2015,hayne2017}. Insolation, albedo,
emissivity, and the upper-$\zborestem$ $K(T,z)$ are therefore correctly
captured by both models, and the
meter-scale-sensor discrepancy in
Table~\ref{tab:headline} is not a surface-driven artifact.

At A17 the Hayne site fit runs surface-warmer than Diviner by
$+4.1$~K
while the Martinez forward sits at $-0.8$~K. The mechanism is direct: the
high meter-scale $K_d^{*}=7.08$ couples
more heat upward
through the diurnal skin layer than the globally calibrated $K(T,z)$
assumes, raising the dayside response. The per-site Hayne fit absorbs this
penalty at the surface to fit the
meter-scale temperature gradient; the surface-calibrated
Martinez form remains essentially unbiased at the surface and pays a
factor-of-$\sim$1.6 penalty in the meter-scale-sensor RMSE relative to
the Hayne site fit (Table~\ref{tab:headline}). Within the Hayne $K(T,z)$ family,
no single parameter choice simultaneously optimises the
surface-brightness constraint of \citet{williams2017diviner} and the
meter-scale-gradient constraint of \citet{nagihara2018} at A17; the
$K(T,\rho)$ family narrows but does not close that gap. The inter-site
difference flagged by the
meter-scale retrieval reappears in the surface-temperature
data, in the form of which model fits which site at the surface and by how
much.

\begin{figure}[!htbp]
\centering
\includegraphics[width=\linewidth]{figures/fig_diviner_closure.pdf}
\caption{Surface-$T$ closure against the \citet{williams2017diviner} Diviner
GCP diurnal composite (gray markers). Teal solid: Hayne form at the per-site
retrieved $K_d^{*}$; violet dashed: \citet{martinez2021} at published
coefficients. Local-solar-time axis centered on midnight (shaded band LST
18--06 = night); full-cycle RMSE\,/\,mean bias per panel in the inset.}
\label{fig:diviner}
\end{figure}

\subsection{Physical plausibility of the retrieved $K_d^{*}$}%
\label{sec:disc-plausibility}

\paragraph*{Apollo 15.} The retrieval gives
$K_d^{\mathrm{A15}}=4.60\,\si{\milli\watt\per\meter\per\kelvin}$. This value
sits between the global \citet{hayne2017} value (3.4) and the regional
\citet{vasavada2012} orbital estimate ($\sim$7), both parameters of the same contact family; at matching
definition, its implied effective conductivity at the sensors
(9.3 mW m$^{-1}$ K$^{-1}$) reproduces the original \citet{langseth1976}
in-situ reduction ($10\pm1$) to within $\sim$7\%. Porosity provides the physical mechanism: the effective-medium
scaling $K_\mathrm{eff}\propto(1-\phi)^2$ of \citet{wood2020}, applied
across the Apollo-core porosity envelope $\phi\!\sim\!0.30$--0.45
\citep{mitchell1973,carrier1991} and combined with the radiative term of
the Hayne form, brackets $\sim$3--6 mW m$^{-1}$ K$^{-1}$. The retrieved
4.60 falls inside that range. The order-of-magnitude gap between the in-situ effective value and direct
laboratory measurements \citep{cremers1971,horai1981} reflects the standard
lab-vs-in-situ systematic: laboratory samples lose the meter-scale
grain-contact network and the radiative pathways that dominate the in-situ
response. A15 is physically admissible.

\paragraph*{Apollo 17.} The retrieval gives
$K_d^{\mathrm{A17}}=7.08\,\si{\milli\watt\per\meter\per\kelvin}$,
roughly twice the global Hayne value and 1.5 times
the A15 retrieval. Two
mechanisms acting together cover this
range comfortably. First, solid lunar basalt
is $\sim$30\% more conductive than highland anorthosite \citep{horai1972};
A17 lies in a basaltic embayment of Mare Serenitatis while A15 occupies a
mare/highland margin \citep{korotev2005,wieczorek2006,crawford2015}.
Second, the porosity-driven $K$ scaling at the upper end of the Apollo-core
range contributes a further $\sim$1.5$\times$ enhancement \citep{wood2020}.
The product, $\sim$1.9$\times$, comfortably covers the
retrieved $\sim$1.5$\times$ ratio within the bootstrap uncertainty. Neither contributor is
required at its extreme, and the value sits just above the
\citet{vasavada2012} orbital estimate of $\sim$7.

A complementary test against the \citet{martinez2021}
$K(T,\rho)$ form (\S\ref{sec:disc-prior}, Fig.~S1, Supporting Information)
is consistent: the density that form requires at A17,
$\rho_d^{*}\approx 1877$ kg m$^{-3}$, lies inside the Apollo-core
envelope, independently corroborating that the A17 meter-scale column
is unremarkable in density and that the enhanced conductivity is
carried by composition and contact structure rather than by anomalous
compaction.

The A17 absolute value is conditional on the adopted basal heat flux. The
\citet{saito2007,nagihara2018} reanalyses revise $Q_b^{\rm A17}$ downward
by up to $\sim$30\%; on the direct map
(\S\ref{sec:disc-degeneracy}) a lower assumed flux moves the retrieved
$K_d^{*}$ up, to $\sim$10.1 mW m$^{-1}$ K$^{-1}$ (and fits the
record slightly better), while the upper bound lowers it to $\sim$6.8.
The flux debate therefore bounds $K_d^{*}$(A17) between $\sim$6.8 and
$\sim$10.1 --- all well above the global value. The
factor-of-$\sim$2 excess over the global Hayne
value therefore cannot be read
as a measured absolute conductivity independent of the $Q_b$ assumption.

A site--depth ambiguity also deserves explicit mention: the two
sensor sets barely overlap in depth (80--139 cm at A15, 130--234 cm at
A17), so a conductivity that continues to increase below
$\sim$1.5 m --- continued compaction rather than a lateral site
difference --- could in principle mimic the contrast. Two internal
diagnostics limit this alternative. Within the A17 set itself, which
spans 1.1 m, a depth-constant $K_d$ leaves only the 0.40 K residual
RMSE of Table~\ref{tab:headline} and predicts the withheld deepest
sensor (234 cm) to 0.25 K (\S\ref{sec:consistency}); a $K(z)$ still
rising across that span would curve the profile detectably. The
borestem-cut sweep similarly leaves both retrievals unchanged as the
A15 set is restricted to its deepest five sensors ($z_b=90$ cm,
\S\ref{sec:results-kd}). Depth-dependent compaction below the probed
range cannot be excluded with two boreholes, but it is not required by
the data, whereas the compositional difference between the sites is
independently established.

A17 remains the more conductive site at the
published basal heat fluxes and across moderate revisions; only the
most adverse corner of the independent-$Q_b$ envelope
(\S\ref{sec:disc-degeneracy}) erases the contrast. Within every
admissible combination of composition and porosity the sign is
unchanged, even where the absolute A17 value is not.

\FloatBarrier
\subsection{Limitations and caveats}\label{sec:limitations}
Six caveats accompany any use of the result. \emph{(i)} With only two
in-situ boreholes, the inter-site difference is insufficient to establish
regional heterogeneity across the lunar surface; distinguishing genuine
regional variability from local heterogeneity at two specific borehole
locations requires additional landed thermometry. \emph{(ii)} The
meter-scale temperature gradient constrains the ratio
$Q_b/K_d$, not $K_d$ alone (\S\ref{sec:disc-degeneracy}); the data favor
but do not establish a conductivity contrast over a uniform conductivity
with differing basal flux. \emph{(iii)} The contrast magnitude
is not resolved at the 95\% level even at the published
basal heat fluxes ($\Delta K_d^{*}=2.3$, bootstrap CI $[-0.1,3.6]$
spanning zero, $P_{\rm boot}{\approx}0.03$; $\sim$97\% of bootstrap draws favor the
A17 ordering). Its \emph{ordering}, however, is robust: on the directly
solved RMSE$(K_d,Q_b)$ surface the joint posterior keeps
$P(K_d^{\rm A17}>K_d^{\rm A15})\ge 99\%$ across the
\citet{saito2007,nagihara2018} $Q_b$ reanalysis envelope, and even the
most adverse independent $Q_b$ pairing only reduces the contrast to
$\sim$$+1.5$ mW m$^{-1}$ K$^{-1}$ without erasing it. What the flux
debate does forfeit is the \emph{absolute} magnitude of both $K_d$
values and of their contrast. These two robustness statements ---
the $\sim$97\% sensor bootstrap and the $\ge$99\% $Q_b$-envelope
ordering --- address orthogonal uncertainty axes and are not combined
into a single joint probability; a hierarchical treatment carrying
both simultaneously (caveat vii) would leave the sign intact but widen
the interval. \emph{(iv)} The
retrieval adopts the
\citet{hayne2017} functional form for $K(T,z)$; alternative piecewise
parameterizations are not exhausted, and a full per-site
recalibration of
the \citet{martinez2021} coefficients (whose density lever is already
admissible at both sites but saturates at A17,
\S\ref{sec:disc-prior}) is the immediate next step. \emph{(v)} The
diurnal-amplitude depth profile (Fig.~\ref{fig:amplitude}) serves only as a
borestem diagnostic and not as a frequency-domain validation, because below
$\zborestem$ both models predict amplitudes well below the 1971--1977
digitization noise floor. \emph{(vi)} The insolation model is the
synodic cosine forcing of \S\ref{sec:methods} with no local topography.
This is least realistic at Apollo~17, whose probes sit on the
Taurus--Littrow valley floor between massifs that shadow the site near
the terminators and re-radiate toward it at midday
\citep{langseth1976}; terrain effects of this kind plausibly contribute
to the $+4.1$ K surface-temperature bias in the Diviner comparison
(\S\ref{sec:consistency}). The retrieval fits only the
sub-skin \emph{lunation-mean} temperatures, so terrain enters through its
lunation-mean perturbation; that perturbation is not negligible a
priori --- horizon blocking removes of order 1\% of the lunation-mean
insolation at a valley-floor site, the same order as the albedo
envelope of Table~\ref{tab:errorbudget} --- but it is partially offset
by terrain infrared emission and reflected sunlight, which return flux
of the same order with the opposite sign. A terrain-resolved forcing
carrying both channels is the natural refinement for any future
re-analysis. \emph{(vii)} The stability windows end at different
epochs across sensors while the record is still warming at depth
\citep{nagihara2018}, so the per-sensor $T_\mathrm{eq}$ values mix
epochs; the bootstrap resamples sensors as exchangeable, ignoring
probe-level calibration offsets and common temporal drift. The
common-epoch sensitivity is now quantified
(\S\ref{sec:consistency}): A15 is unaffected, while A17 moves by
$-0.89$ mW m$^{-1}$ K$^{-1}$ (carried as $\sigma_\mathrm{epoch}$ in
Table~\ref{tab:errorbudget}) and the contrast becomes more marginal
($P_{\rm boot}=0.04$) while retaining its sign. A probe-clustered
(hierarchical) uncertainty treatment remains follow-on work.

% =====================================================================
\FloatBarrier
\section{Conclusions}\label{sec:conclusions}

The two Apollo boreholes require distinctly different
meter-scale conductivities. Holding the \citet{hayne2017}
functional form fixed and sweeping $K_d$ alone against the
meter-scale-sensor RMSE, the per-site retrieval
yields
$K_{d,\mathrm{A15}}^{*}=4.60^{+0.86}_{-0.21}$ and
$K_{d,\mathrm{A17}}^{*}=7.08^{+0.53}_{-0.52}\,\si{\milli\watt\per\meter\per\kelvin}$
($1\sigma$ non-parametric bootstrap, $N_\mathrm{boot}=1500$) at the
published basal heat fluxes, with the inter-site contrast
only marginally resolved ($\Delta K_d^{*}=2.3$, 95\% CI
$[-0.1,3.6]$ spanning zero, $P_{\rm boot}{\approx}0.03$). The
meter-scale-sensor RMSE decreases from
1.09 K and 0.89 K under the published global
$K_d=3.4\,\si{\milli\watt\per\meter\per\kelvin}$ to 1.00 K and 0.40 K under
the per-site retrieval; the \citet{martinez2021} $K(T,\rho)$ forward
evaluation, calibrated against the same family of orbital products,
reproduces both sites to within $\sim$1 K without per-site
freedom, though it cannot match the sharpened per-site A17 fit
(Table~\ref{tab:headline}). The inter-site requirement is robust to the
borestem exclusion depth, the stability-window threshold, a joint
$K_d$--$H$ fit, and the independent Diviner surface-temperature closure
check.

The absolute $K_d^{*}$ values are conditional on the adopted $Q_b$. Because
the meter-scale temperature gradient constrains the ratio
$Q_b/K_d$ rather than $K_d$ alone, the inter-site difference admits two
interpretations: a conductivity contrast or a basal-flux contrast. A pooled
three-model comparison cannot separate the two readings
($\Delta\mathrm{AICc}\le 0.5$ across all three models); the per-site
evidence at A17 is what is decisive,
$\Delta\mathrm{AICc}\approx 23$ for the site fit over the global value
(Table~\ref{tab:modelcomp}) without
excluding the alternative. The quantity constrained tightly at each
borehole is the meter-scale temperature gradient itself.

The immediate application is to forward thermal models used by upcoming
lunar sub-surface RTM retrievals. These retrievals require a
meter-scale $T(z)$ boundary condition that orbital-only
conductivity models cannot supply on a per-site basis. The per-site
$K_d^{*}$ values reported here, evaluated in the Hayne thermal column,
provide that boundary condition (\S\ref{sec:intro-tsukimi}); the upcoming
TSUKIMI terahertz mission \citep{tsukimi2022} is one immediate beneficiary.
A full per-site recalibration of the
\citet{martinez2021} form, and the
extension of the meter-scale $T(z)$ column to include
surface-shadowing effects on the RTM side, are both deferred to follow-on
work. Extending the analysis beyond two boreholes, to a regionally resolved
treatment of the meter-scale lunar thermal state, is
fundamentally a measurement problem that requires
meter-scale temperature data from additional
landed instruments (e.g.\ Chang'E and ChaSTE-class experiments).

% =====================================================================
%  AGU back matter: Acknowledgments -> Open Research -> Conflict of
%  Interest -> References (the order required by the AGU template).
% =====================================================================
\section*{Acknowledgments}
The authors thank the original Apollo HFE Principal Investigators and the
\citet{nagihara2018} team for restoring the 1975--1977 portion of the
1971--1977 record from the original tapes (the 1971--1975 portion was
retained from the original PI processing), without which this work would not
have been possible.


% =====================================================================
\section*{Open Research}\label{sec:openresearch}
\paragraph*{Availability Statement.}
The restored Apollo~15 and~17 Heat-Flow Experiment temperature record
(1971--1977) analyzed here is openly available as the supporting
information of \citet{nagihara2018} (\url{https://doi.org/10.1029/2018JE005579}).
The Diviner Lunar Radiometer Global Cumulative Product used for the
out-of-sample surface-temperature closure \citep{williams2017diviner} is
available from the NASA Planetary Data System Geosciences Node
(\url{https://pds-geosciences.wustl.edu}; LRO-L-DLRE-5-GCP-V1.0,
\url{https://doi.org/10.17189/1520650}). The
repository bundles the NASA SPICE Toolkit kernels (JPL DE440;
\citealp{acton1996,acton2018}) used by its ephemeris module. The 1-D
Crank--Nicolson heat solver, the
ephemeris-driven solar-forcing routines, the $K_d$ retrieval and bootstrap
code, the validation diagnostics, and the data and scripts needed to
reproduce every figure and table in this manuscript are archived at
\url{https://github.com/rp3gregorio/Lunar-HFE}; a
Zenodo-archived release with a citable DOI will be issued at the time of
acceptance and added to the reference list.


% =====================================================================
\section*{Conflict of Interest Disclosure}
The authors declare there are no conflicts of interest for this manuscript.

% =====================================================================
\section*{Nomenclature}
\label{sec:nomenclature}
{\small
\setlength{\tabcolsep}{4pt}
\renewcommand{\arraystretch}{1.0}

\noindent\textit{Symbols.}\par\smallskip\noindent
\begin{tabular}{@{}p{0.10\linewidth}p{0.40\linewidth}p{0.28\linewidth}p{0.16\linewidth}@{}}
\toprule
\textbf{Symbol} & \textbf{Definition} & \textbf{Value / units} & \textbf{Source} \\
\midrule
$K_d$              & deep asymptote of the contact $K_c$ (retrieved); effective $K$ at sensors ${\approx}2\times$        & 4.60, 7.08 mW\,m$^{-1}$\,K$^{-1}$         & this work \\
$K_s$              & surface regolith thermal conductivity                  & $7.4\times10^{-4}$ W\,m$^{-1}$\,K$^{-1}$  & \citet{hayne2017} \\
$H$                & e-folding depth of $K(z)$ transition                  & 0.06 m                                       & \citet{hayne2017} \\
$\chi$             & radiative conductivity multiplier coefficient         & 2.7                                          & \citet{hayne2017} \\
$T_\mathrm{ref}$   & normalization temperature for radiative term          & 350 K                                        & \citet{hayne2017} \\
$\rho_s,\rho_d$    & surface / asymptotic regolith bulk density                  & 1100, 1800 kg\,m$^{-3}$                      & \citet{hayne2017} \\
$c_p(T)$           & specific heat (4th-order polynomial in $T$)           & J\,kg$^{-1}$\,K$^{-1}$                       & \citet{hemingway1973,hayne2017} \\
$Q_b$              & basal heat flux                                       & 21 (A15), 16 (A17) mW\,m$^{-2}$              & \citet{langseth1976,nagihara2018} \\
$S_0$              & solar constant (irradiance at 1 AU)                   & 1361 W\,m$^{-2}$                             & \citet{kopp2011} \\
$A$, $\varepsilon$ & broadband Bond albedo, emissivity                     & site-specific;\ 0.95                         & \citet{vasavada2012,bandfield2015} \\
$\sigma$           & Stefan--Boltzmann constant                            & $5.6704\!\times\!10^{-8}$ W\,m$^{-2}$\,K$^{-4}$ & CODATA 2018 \\
$\theta_\odot(t)$  & solar zenith angle at site                            & rad                                          & synodic approximation, \S\ref{sec:methods} \\
$F_\odot(t)$       & incident solar irradiance at surface                  & W\,m$^{-2}$                                  & this work, eq.~\ref{eq:insolation} \\
$T(z,t),\,T_s$     & subsurface and skin ($z\!=\!0$) temperature           & K                                            & this work \\
$z$                & depth coordinate (positive downward)                  & m                                            & --- \\
$\zborestem$       & borestem-exclusion depth (\emph{a priori} cut)        & 80 cm                                        & this work, \S\ref{sec:borestem} \\
$\delta$           & diurnal thermal skin depth                            & $\sim$3--5 cm                                & \citet{hayne2017} \\
$\Delta K_d$       & inter-site contrast $K_d^{\rm A17}\!-\!K_d^{\rm A15}$ & 2.31 mW\,m$^{-1}$\,K$^{-1}$                  & this work \\
RMSE               & meter-scale-sensor root-mean-square residual                 & K                                            & this work \\
$N$                & number of meter-scale sensors                                & 7 (A15), 16 (A17)                            & \citet{nagihara2018} \\
$N_\mathrm{boot}$  & non-parametric bootstrap resamples                    & 1500                                         & \citet{efron1993} \\
$P_{\rm boot}$     & bootstrap tail proportion $P(\Delta K_d^{*}\le 0)$ (not a null-centered $p$-value) & 0.031                                         & this work \\
\bottomrule
\end{tabular}

\medskip
\noindent\begin{minipage}[t]{0.49\linewidth}
\textit{Subscripts.}\par\smallskip
\begin{tabular}{@{}ll@{}}
\toprule
$s$            & surface (regolith, skin) \\
$d$            & asymptotic (deep) regolith  \\
$b$            & basal (lower BC) \\
$\mathrm{ref}$ & reference (normalization) \\
$\odot$        & solar \\
$\mathrm{eq}$  & stability-window mean \\
$\mathrm{boot}$ & bootstrap quantity \\
\bottomrule
\end{tabular}
\end{minipage}\hfill
\begin{minipage}[t]{0.49\linewidth}
\textit{Superscripts.}\par\smallskip
\begin{tabular}{@{}ll@{}}
\toprule
$^{*}$            & best-fit / retrieved (e.g.\ $K_d^{*}$) \\
$^{\rm A15,A17}$  & per-site qualifier \\
$^{\rm ref}$      & reference value \\
$^{\rm HFE}$      & Apollo HFE-measured \\
$\langle\cdot\rangle$ & stability-window time average \\
\bottomrule
\end{tabular}
\end{minipage}

\medskip
\noindent\textit{Acronyms.}\par\smallskip\noindent
\begin{tabular}{@{}ll@{}}
\toprule
HFE          & Heat-Flow Experiment (Apollo~15 and~17 buried thermometers) \\
TG / TR      & gradient-bridge / ring-bridge sensor (HFE sensor types) \\
LST          & local solar time \\
Diviner GCP  & Diviner Lunar Radiometer Global Cumulative Product \\
MCMC         & Markov-chain Monte Carlo \\
BCa          & bias-corrected accelerated bootstrap \\
SPICE        & NASA ancillary information system for planetary geometry \\
DE440        & JPL planetary and lunar ephemeris release 440 \\
RTM          & radiative-transfer model \\
TSUKIMI      & lunar mission consuming the meter-scale $T(z)$ retrieved here \\
\bottomrule
\end{tabular}
}

\bibliography{references}

\end{document}
